\chapter{Competition creates competition}\label{ch:ccc}

\begin{chapterplan}
The question of this chapter: why can a stronger incumbent raise the probability that a challenger prepares,
when it never raises what preparation earns at any belief? After this chapter you can do five things.
\begin{itemize}
\item Tell the story of three economies with a weak, a strong and a very strong incumbent, and say what the
  price does in each.
\item State Proposition~2 with every hypothesis, and say what each of the conditions (H1), (H2) and (H3) compares.
\item Derive when some trading cost satisfies (H2), and say why a prior away from one half makes this harder.
\item Rebuild the proof economy by economy: orders first, then posteriors, then entry, then prices, with
  uniqueness under all four tie rules.
\item Derive entry, high-class ownership and the takeover-share condition against the strong incumbent, and use
  the decisive-advantage lemma to obtain the case in which class \(H\) outbids the incumbent.
\end{itemize}
\end{chapterplan}

\section{Three economies in words}\label{sec:ccc-story}

At fixed information, a stronger incumbent hurts a challenger. Recall from
Chapter~\ref{ch:auction} that a strengthening of the incumbent never raises the gross profit \(g_H\) of class
\(H\) or the gross profit \(g_L\) of class \(L\) (Proposition~1(a)). So gross profit \(B(\mu)\) at every fixed
belief \(\mu\) never rises. The auction-entry literature calls this deterrence; the paper cites Fishman (1988)
and Hirshleifer and Png (1989) for it. If the challenger's belief stayed the same, a stronger incumbent could not
raise entry.

The paper adds a second force with the same cause. A stronger incumbent never narrows the target-payoff spread
\(\DT=t_H-t_L\) (Proposition~1(b)). The spread is what the informed investor trades on. So a stronger incumbent can
make informed trade pay, and informed trade puts information into the price. The challenger reads the price
before it decides whether to prepare.

The main result compares three economies. They differ only in the incumbent's value distribution: a weak
incumbent \(F_0\), a strong incumbent \(F_1\) that strengthens \(F_0\), and a very strong incumbent \(F_2\) that
strengthens \(F_1\). The classes, the reserve, the prior, the order slot and the costs are the same in all three.
\begin{itemize}
\item \emph{The weak incumbent.} The spread is narrow. An informed order earns at most the spread, and the
  spread is below the trading cost \(k\). So the investor never trades, and the price is silent: after every
  order the belief stays at the prior \(\pi\). At the prior, gross profit is below the high cost \(c_H\). Only a
  zero-cost challenger prepares, and entry equals the floor \(\rho\).
\item \emph{The strong incumbent.} The spread is wide. The investor's residual advantage covers \(k\) under
  every strategy, so class \(H\) buys and class \(L\) sells. A purchase lifts the belief to the top posterior
  \(M(\pi)\). At \(M(\pi)\), gross profit exceeds \(c_H\), so a purchase recruits a high-cost challenger. Entry
  rises above \(\rho\).
\item \emph{The very strong incumbent.} The spread is at least as wide, so the investor still trades and the
  price still speaks. But gross profit is so low that it stays below \(c_H\) even at \(M(\pi)\). No price recruits
  a high-cost challenger, and entry returns to \(\rho\).
\end{itemize}
Entry first rises and then falls. Figure~\ref{fig:ccc-sketch} draws this pattern.

\begin{figure}[ht]
\centering
\begin{tikzpicture}[x=1cm,y=1cm,font=\small]
\draw[-{Stealth}] (0,0) -- (10.8,0) node[right]{incumbent strength};
\draw[-{Stealth}] (0,0) -- (0,4.8) node[above]{entry};
\draw[dashed,rulegrey] (0,1.2) -- (10.4,1.2);
\node[left] at (0,1.2) {\(\rho\)};
\draw[dotted,rulegrey] (0,3.5) -- (5.2,3.5);
\node[left] at (0,3.5) {\(\Ent_1\)};
\draw (-0.08,4.3) -- (0.08,4.3);
\node[left] at (0,4.3) {\(1\)};
\foreach \x/\lab in {2/F_0, 5.2/F_1, 8.4/F_2} {
  \draw (\x,-0.08) -- (\x,0.08);
  \node[below] at (\x,-0.05) {\(\lab\)};
}
\fill[accent] (2,1.2) circle (2.6pt);
\fill[accent] (5.2,3.5) circle (2.6pt);
\fill[accent] (8.4,1.2) circle (2.6pt);
\draw[-{Stealth},accent] (2.35,1.5) -- (4.85,3.2) node[midway,above left]{rises};
\draw[-{Stealth},accent] (5.55,3.2) -- (8.05,1.5) node[midway,above right]{falls};
\node[below,align=center,text=notegrey] at (2,-0.55) {price silent};
\node[below,align=center,text=notegrey] at (5.2,-0.55) {price speaks,\\a purchase recruits};
\node[below,align=center,text=notegrey] at (8.4,-0.55) {price speaks,\\no price recruits};
\end{tikzpicture}
\caption{A sketch of Proposition~2. Each dot is the unique equilibrium entry of one economy. The dashed line is
the entry floor \(\rho\). The proposition compares only these three economies; Chapter~\ref{ch:coexistence}
describes the strengths between them. The sketch has no scale.}
\label{fig:ccc-sketch}
\end{figure}

\begin{intuition}[two forces with one cause]
A stronger incumbent moves part of the class value gap away from the challenger and toward the target. The
challenger keeps no more than before, so its gross profit does not rise at any belief. The target's stock receives
at least as much of the gap, so the stock becomes at least as sensitive to the challenger's class. From \(F_0\) to \(F_1\) the second effect switches information
on: the belief after a purchase jumps from \(\pi\) to \(M(\pi)\), and this jump outweighs the fall in profit at
each belief. From \(F_1\) to \(F_2\) information is already on, so only the fall in profit acts.
\end{intuition}

\begin{analogy}[a loudspeaker with a fee]
Think of the informed investor as a loudspeaker that costs \(k\) to switch on. It speaks only when its message is
worth more than the fee. The worth of the message grows with the spread, so a stronger incumbent makes speech
more likely. The challenger listens, and it acts on good news when good news makes preparation pay.

The analogy breaks in three places. First, the investor does not want to inform anyone; it trades for its own
profit, and the price carries the message as a side effect. Second, the message is noisy, because a liquidity
trader can place the same order. Third, the worth of the message depends on how the listener reacts, and market
makers put that reaction into the price.
\end{analogy}

\section{The objects of the result}\label{sec:ccc-setup}

The proposition uses objects from four earlier chapters. This section recalls each one in words, with its
formula. Keep this page open while you read the proof.

\subsection{The auction}

\begin{itemize}
\item The challenger belongs to class \(H\) or class \(L\). The prior is \(\pi=\Pr(j=H)\), with \(0<\pi<1\). Given
  its class \(j\), the challenger's value \(\theta\) has the distribution \(G_j\) with a finite mean. Write
  \(D(y)=G_L(y)-G_H(y)\).
\item The class ranking (equation (1) of the paper) says \(G_H(y)\le G_L(y)\) for every \(y\), and
  \(G_H(p^-)=G_L(p^-)=0\). So \(D(y)\ge0\) everywhere, and every value meets the reserve \(p>0\).
\item The incumbent's value \(R\) has a continuous distribution \(F\). A strengthening of \(F\) is a continuous
  distribution \(\tilde F\) with \(\tilde F\le F\) at every point (Section~\ref{sec:toolkit-fosd}).
\item Without entry the target receives \(t_0=p[1-F(p)]\) in expectation. After entry by a class-\(j\) challenger
  it receives \(t_j\), and the challenger expects the gross profit \(g_j\). The spread is \(\DT=t_H-t_L\).
\item Write \(w_j=t_j-t_0\) for the gain in expected proceeds that entry brings in class \(j\), and
  \(w(\mu)=(1-\mu)w_L+\mu w_H=t_L-t_0+\DT\mu\). An admissible challenger raises proceeds to at least the reserve,
  so \(t_j\ge t_0\) and \(w_H,w_L\ge0\).
\item Gross profit at the belief \(\mu=\Pr(j=H)\) is \(B(\mu)=g_L+\mu(g_H-g_L)\).
\item A prepared class-\(j\) challenger beats the incumbent with probability
  \(W_j(F)=\E_j\bigl[F(\theta)\1\{\theta\ge p\}\bigr]\) (Chapter~\ref{ch:model}). Under the class ranking every value
  meets the reserve, so \(W_j(F)=\E_j[F(\theta)]\). Class \(H\) \emph{outbids} the incumbent \(F\) when
  \(W_H(F)=1\).
\end{itemize}

Proposition~1 (Chapter~\ref{ch:auction}) gives the auction objects as integrals:
\begin{equation}\label{eq:ccc-prop1}
\DT(F)=pF(p)D(p^-)+\int_p^\infty[1-F(y)]\,D(y)\,dy,\qquad
g_j(F)=\int_p^\infty[1-G_j(y)]\,F(y)\,dy.
\end{equation}
Part (a): a strengthening weakly lowers \(g_H\) and \(g_L\), and so gross profit at every fixed belief. Part (b):
under the class ranking, a strengthening weakly raises \(\DT\). Both changes are weak in general; each is strict
when its integral of the difference \(F-\tilde F\) is positive.

Two facts follow from \eqref{eq:ccc-prop1} under the class ranking. The proof uses both.

\begin{derivation}[the spread is never negative and the profit gap is never negative]
The class ranking gives \(D(p^-)=0\) and \(D(y)\ge0\). So
\begin{align*}
\DT(F)&=0+\int_p^\infty[1-F(y)]\,D(y)\,dy \why{\(D(p^-)=0\)}\\
&\ge0 \why{\(1-F\ge0\) and \(D\ge0\)}.
\end{align*}
For the profit gap, subtract the two lines of \eqref{eq:ccc-prop1}:
\begin{align*}
g_H(F)-g_L(F)&=\int_p^\infty[1-G_H(y)]\,F(y)\,dy-\int_p^\infty[1-G_L(y)]\,F(y)\,dy \why{\eqref{eq:ccc-prop1}}\\
&=\int_p^\infty[G_L(y)-G_H(y)]\,F(y)\,dy \why{one integral}\\
&=\int_p^\infty D(y)\,F(y)\,dy \why{definition of \(D\)}\\
&\ge0 \why{\(D\ge0\) and \(F\ge0\)}.
\end{align*}
Gross profit is linear in the belief. For two beliefs \(\mu\) and \(\mu'\),
\begin{align*}
B(\mu)-B(\mu')&=[g_L+\mu(g_H-g_L)]-[g_L+\mu'(g_H-g_L)] \why{definition of \(B\)}\\
&=(\mu-\mu')(g_H-g_L) \why{\(g_L\) cancels}.
\end{align*}
So \(B\) never falls when the belief rises, and it rises strictly when \(g_H>g_L\).
\end{derivation}

\subsection{Preparation and the order slot}

\begin{itemize}
\item After it sees the price, the challenger learns its preparation cost \(C\in\{0,c_H\}\), with
  \(\Pr(C=0)=\rho\), \(0<\rho<1\), and \(c_H>0\). The cost is independent of the class, the values and the slot.
  A zero-cost challenger prepares after every price.
\item Write \(\zeta(\bar P)\) for the probability that a high-cost challenger prepares after the price \(\bar P\).
  Entry after the order \(x\) is \(e(x)=\rho+(1-\rho)\,\zeta(P(x))\), so \(e(x)\ge\rho\). This is the entry
  floor.
\item In the order slot, the informed investor arrives with probability \(\lambda\in(0,1)\). It sees the class and
  places \(q\in\{-1,0,1\}\); a nonzero order costs \(k>0\). Otherwise a liquidity trader sells, abstains or buys with
  probabilities \(\beta\), \(1-2\beta\) and \(\beta\).
\item A liquidity order in one direction has probability \(b=(1-\lambda)\beta\). A liquidity abstention has
  probability \(b_0=(1-\lambda)(1-2\beta)\). The mass ratio is \(u=\lambda/b>0\).
\item Write \(\sigma_j(x)\) for the probability that the investor in class \(j\) places the order \(x\). The order
  \(x\) then has probability \(f_j(x)=(1-\lambda)\eta(x)+\lambda\sigma_j(x)\) in class \(j\), where \(\eta\) gives
  the liquidity probabilities \((\beta,1-2\beta,\beta)\) of the orders \((-1,0,1)\). Write \(\delta_q\) for the
  strategy that places \(q\) for sure.
\end{itemize}

\subsection{Beliefs, prices and trade}

\begin{itemize}
\item Bayes' rule in odds form (Section~\ref{sec:toolkit-odds}) gives the posterior \(\mu(x)\) after the order \(x\):
  \[
  \frac{\mu(x)}{1-\mu(x)}=\frac{\pi}{1-\pi}\cdot\frac{f_H(x)}{f_L(x)}.
  \]
\item The posterior bounds \eqref{eq:beliefs-bounds} are
  \[
  M(\pi)=\frac{\pi(1+u)}{\pi(1+u)+1-\pi},\qquad m(\pi)=\frac{\pi}{\pi+(1-\pi)(1+u)}.
  \]
  After a nonzero order, \(m(\pi)\le\mu(x)\le M(\pi)\) under every investor strategy and for every \(\beta<1/2\)
  (Lemma~A.1(b)). The same bound after the zero order needs \(\beta\le1/3\) (Lemma~A.1(c)). When class \(H\) buys
  and class \(L\) sells for sure, the posteriors after \((-1,0,1)\) are \((m(\pi),\pi,M(\pi))\) (Lemma~A.1(d)).
\item The challenger's belief at a price averages the posteriors of the orders that share this price. So it lies
  between the smallest and the largest of these posteriors.
\item Write \(\kappa(\pi)=\min\{m(\pi),1-M(\pi)\}\).
\item When \(g_H>g_L\), the preparation threshold \eqref{eq:beliefs-threshold} is \(\tau=(c_H-g_L)/(g_H-g_L)\), and
  \(B(\mu)\ge c_H\) exactly when \(\mu\ge\tau\).
\item The exclusion bound (Section~\ref{sec:beliefs-exclusion}, Corollary~A.1): if \(0<\beta\le1/3\),
  \(g_H\ge g_L\) and \(B(M(\pi))<c_H\), then no high-cost challenger prepares after any order.
\item The price formula \eqref{eq:beliefs-price}: in every candidate,
  \[
  P(x)=t_0+e(x)\bigl[t_L-t_0+\DT\,\mu(x)\bigr]=t_0+e(x)\,w(\mu(x)).
  \]
  It uses only Bayes' rule, competitive prices, truthful bids and the independence of the cost.
\item The price reveals the posterior (Section~\ref{sec:beliefs-reveal}): if \(\DT>0\) and \(\rho>0\), two orders
  with one price have one posterior, and the challenger's belief at \(P(x)\) is \(\mu(x)\). This is Lemma~A.2(c).
  The lemma also assumes \(g_H>g_L\), but the proof of part (c) uses only the price formula, \(\DT>0\) and
  \(e\ge\rho>0\).
\item The trading returns \eqref{eq:trading-returns} of the two correct orders are
  \(K_H=e(1)\,\DT\,[1-\mu(1)]\) and \(K_L=e(-1)\,\DT\,\mu(-1)\). Class \(H\) earns \(U_H(1)=K_H-k\) from a purchase,
  class \(L\) earns \(U_L(-1)=K_L-k\) from a sale, and abstention earns \(U_j(0)=0\). A wrong-signed order earns at
  most \(-k\) (Lemma~A.3).
\item The trading bound \eqref{eq:trading-bound}: under every candidate profile,
  \[
  \rho\,\kappa(\pi)\,\DT\le\rho\,[1-M(\pi)]\,\DT\le K_H\le\DT,\qquad
  \rho\,\kappa(\pi)\,\DT\le\rho\,m(\pi)\,\DT\le K_L\le\DT.
  \]
\end{itemize}

\subsection{Equilibrium, candidates and tie rules}

An equilibrium (Chapter~\ref{ch:model}, Definition~A.1) lists the investor's order distributions, the posteriors,
the prices, the challenger's belief at each price and the high-cost preparation rule, with truthful bids. It
meets six conditions: (E1) Bayes' rule after each order; (E2) Bayes' rule at each price; (E3) optimal high-cost
preparation; (E4) competitive prices; (E5) optimal orders; (E6) truthful bids.

A \emph{candidate} meets (E1) to (E4) and (E6) for some pair of order distributions. An equilibrium is a
candidate that also meets (E5). The uniqueness proof derives bounds that hold for every candidate. So it never
assumes the tuple that it wants to find.

Two tie rules apply. An indifferent investor may mix over its best orders, or it abstains. An indifferent
high-cost challenger prepares, or it may prepare with any probability. The four pairs are the four tie rules.
The default tie rules are ``may mix'' and ``prepares''.

For the three economies, write \(\Delta_i=\DT(F_i)\), \(B_i=B_{F_i}\), and \(t_{0,i}\), \(t_{L,i}\), \(g_{j,i}\),
\(w_i(\mu)\) and \(\tau_i\) for the objects against \(F_i\), with \(i\in\{0,1,2\}\). Inside the proof, write
\(M=M(\pi)\), \(m=m(\pi)\) and \(\kappa=\kappa(\pi)\).

\section{The statement}\label{sec:ccc-statement}

Here is the result as the paper states it. The box copies every hypothesis.

\begin{paperresult}{Proposition 2 (Competition creates competition)}
Let the classes satisfy the class ranking and have finite means, let \(F_1\) be a continuous first-order
stochastic strengthening of the continuous incumbent distribution \(F_0\), and let \(F_2\) strengthen \(F_1\) in
the same sense. Fix \(0<\pi,\rho,\lambda<1\), \(0<\beta<1/2\), \(c_H>0\) and \(k>0\). Write \(\Delta_i=\DT(F_i)\),
\(B_i=B_{F_i}\), and \(t_{0,i}\) and \(t_{L,i}\) for the proceeds \(t_0\) and \(t_L\) against \(F_i\), with
\(b=(1-\lambda)\beta\) and \(\kappa(\pi)=\min\{m(\pi),1-M(\pi)\}\). Suppose that
\[
\begin{array}{ll}
\text{(H1), the high-cost window:}& B_0(\pi)<c_H<B_1(M(\pi)),\\
\text{(H2), the trading-cost window:}& \Delta_0<k<\rho\kappa(\pi)\Delta_1,\\
\text{(H3), the exclusion condition:}& B_2(M(\pi))<c_H.
\end{array}
\]
Each economy has a unique equilibrium outcome under all four tie rules.
\begin{enumerate}[label=(\roman*),leftmargin=2em]
\item \emph{The weak incumbent \(F_0\).} Both classes abstain, every posterior equals \(\pi\), entry is \(\rho\),
  and all orders share the price \(t_{0,0}+\rho[t_{L,0}-t_{0,0}+\Delta_0\pi]\). High-class ownership is
  \(\pi\rho W_H(F_0)\).
\item \emph{The strong incumbent \(F_1\).} Class \(H\) buys and class \(L\) sells with probability one. The
  posteriors after \((-1,0,1)\) are \((m(\pi),\pi,M(\pi))\), and entry is \((\rho,\rho,1)\). Prices rise strictly in
  this order. Class entry and total entry are
  \[
  e_H=\rho+(1-\rho)(b+\lambda),\qquad e_L=\rho+(1-\rho)b,\qquad \Ent_1=\rho+(1-\rho)(b+\pi\lambda)>\rho.
  \]
  High-class ownership rises exactly when \(e_HW_H(F_1)>\rho W_H(F_0)\). Class \(H\)'s takeover share rises exactly
  when
  \begin{equation}\label{eq:ccc-share}
  e_HW_H(F_1)W_L(F_0)>e_LW_H(F_0)W_L(F_1).
  \end{equation}
  If class \(H\) outbids the stronger incumbent, \(W_H(F_1)=1\), it also outbids \(F_0\); high-class ownership rises
  from \(\pi\rho\) to \(\pi e_H\), and its takeover share is nondecreasing, strictly exactly when \(W_L(F_0)>0\). If
  \(W_L(F_0)=0\), that share is one against both incumbents.
\item \emph{The very strong incumbent \(F_2\).} The orders and posteriors are those of part~(ii), entry is \(\rho\)
  after every order, and the three prices differ. High-class ownership is \(\pi\rho W_H(F_2)\).
\end{enumerate}
Against each \(F_i\), with \(\mu_i(x)\) and \(e_i(x)\) the posterior and entry after the order \(x\), prices are
\(P_i(x)=t_{0,i}+e_i(x)[t_{L,i}-t_{0,i}+\Delta_i\mu_i(x)]\), the price belief is \(\mu_i(x)\), high-cost preparation
is \([e_i(x)-\rho]/(1-\rho)\), and bids are truthful.
\end{paperresult}

In the paper, the takeover-share condition \eqref{eq:ccc-share} is equation (6). The hypotheses ask for ranked
classes, three continuous incumbent distributions with \(F_2\le F_1\le F_0\) at every point, interior
probabilities, \(0<\beta<1/2\) and positive costs. Section~\ref{sec:ccc-conditions} reads (H1) to (H3). In each
economy exactly one equilibrium outcome exists, whatever the tie rule. Read the rest one piece at a time.
\begin{itemize}
\item \emph{Part (i).} The investor never trades. The price is the same after every order, so it carries no
  information. Only the zero-cost challenger prepares.
\item \emph{Part (ii).} The investor always trades in the correct direction. A purchase moves the belief to the
  top posterior and recruits a high-cost challenger; the other two orders leave entry at the floor. Entry in each
  class and total entry rise above \(\rho\). The ownership and share statements say when the extra entry turns
  into more acquisitions by class \(H\).
\item \emph{Part (iii).} The investor trades as in part (ii), so the price still speaks. No price recruits a
  high-cost challenger, so entry is back at the floor.
\item \emph{The last paragraph.} In each economy, prices follow the price formula, the challenger reads the
  posterior from the price, and the high-cost preparation probability follows from entry.
\end{itemize}

\begin{watchout}
Copy the domain of \(\beta\) exactly: Proposition~2 states \(0<\beta<1/2\), while the model section and the
exclusion bound use \(\beta\le1/3\). The proof fixes orders with the posterior bounds after nonzero orders, which
hold for every \(\beta<1/2\), and then computes the posteriors exactly. The watchout in
Section~\ref{sec:ccc-conditions} on (H3) explains why the zero order never needs \(\beta\le1/3\).
\end{watchout}

\section{The three conditions in words}\label{sec:ccc-conditions}

Each condition is a pair of strict inequalities or a single strict inequality. For each one, this section says
what it compares, why the proof needs it, and what goes wrong without it. Figure~\ref{fig:ccc-profit} shows (H1) and
(H3) in one picture.

\begin{figure}[ht]
\centering
\begin{tikzpicture}[x=1cm,y=1cm,font=\small]
\draw[-{Stealth}] (0,0) -- (8.8,0) node[below]{belief \(\mu\)};
\draw[-{Stealth}] (0,0) -- (0,5.4) node[above]{gross profit};
\draw (8,-0.08) -- (8,0.08); \node[below] at (8,-0.05) {\(1\)};
\node[below] at (0,-0.05) {\(0\)};
\draw[dashed,ink] (0,2.6) -- (8.3,2.6);
\node[left] at (0,2.6) {\(c_H\)};
\foreach \x/\lab in {1.9/m, 3.3/\pi, 5.4/M} {
  \draw[dotted,rulegrey] (\x,0) -- (\x,4.9);
  \node[below] at (\x,-0.05) {\(\lab\)};
}
\draw (4.84,-0.08) -- (4.84,0.08); \node[below,text=notegrey] at (4.84,-0.05) {\(\tau_1\)};
\draw[thick,accent] (0,0.5) -- (8,5.0) node[right]{\(B_0\)};
\draw[thick,ink] (0,0.3) -- (8,4.1) node[right]{\(B_1\)};
\draw[thick,notegrey] (0,0.15) -- (8,3.3) node[right]{\(B_2\)};
\fill[accent] (3.3,2.356) circle (2.2pt);
\fill[ink] (5.4,2.865) circle (2.2pt);
\fill[notegrey] (5.4,2.277) circle (2.2pt);
\end{tikzpicture}
\caption{A picture to keep in mind for (H1) and (H3); the lines show one possible case, with no scale. Each
line is the gross profit \(B_i(\mu)=g_{L,i}+\mu(g_{H,i}-g_{L,i})\) against one incumbent. The dots mark
\(B_0(\pi)<c_H\), \(B_1(M)>c_H\) and \(B_2(M)<c_H\). Each line lies on or below the one before and is no steeper,
because a strengthening never raises \(g_L\) and never raises the gap \(g_H-g_L=\int_p^\infty D(y)F(y)\,dy\).}
\label{fig:ccc-profit}
\end{figure}

\subsection{(H1), the high-cost window}

\emph{What it compares.} (H1) reads \(B_0(\pi)<c_H<B_1(M(\pi))\). It compares the high cost with gross profit at
two beliefs in two economies. On the left is gross profit at the prior against the weak incumbent. On the right is
gross profit at the top posterior against the strong incumbent.

\emph{Why the proof needs it.} The left side gives entry \(\rho\) against \(F_0\): there the price is silent, so the
challenger's belief is \(\pi\), and gross profit at \(\pi\) does not cover \(c_H\). Through Proposition~1(a),
\(B_1(\pi)\le B_0(\pi)<c_H\), so the left side also stops recruitment after the zero order against \(F_1\). The right
side gives recruitment after a purchase against \(F_1\). Together, the two sides force \(g_{H,1}>g_{L,1}\), so the
threshold \(\tau_1\) exists, and \(\pi<\tau_1<M(\pi)\) (Step~0 of the proof).

\emph{What fails without it.} Suppose the left side asked only \(B_1(\pi)<c_H\). Then gross profit at the prior can
exceed \(c_H\) against \(F_0\). When it does, the silent price leaves the belief at \(\pi\), both cost types prepare,
and entry against \(F_0\) is one. Against \(F_1\) entry is \(\Ent_1<1\), so entry falls. Appendix A.3 of the paper
builds such a case from a variant of the example. Suppose instead the right side fails, \(c_H>B_1(M(\pi))\). Orders are still full,
so no price lifts the belief above \(M(\pi)\). No price recruits against \(F_1\), and entry stays at \(\rho\).
Chapter~\ref{ch:example} shows this case in the example.

\begin{checkyourself}
Show that (H1) implies \(B_0(M(\pi))>c_H\). What does this inequality say about the reason why no high-cost
challenger prepares against the weak incumbent?
\label{q:ccc-missing}
\end{checkyourself}

\subsection{(H2), the trading-cost window}

\emph{What it compares.} (H2) reads \(\Delta_0<k<\rho\kappa(\pi)\Delta_1\). It compares the trading cost with the two
ends of the trading bound. On the left is the upper end against the weak incumbent: no correct order earns more
than \(\Delta_0\). On the right is the lower end against the strong incumbent: every correct order earns at least
\(\rho\kappa(\pi)\Delta_1\).

\emph{Why the proof needs it.} The trading bound holds under every candidate. So (H2) fixes the investor's orders
before any order profile is known. Against \(F_0\), every correct order earns at most \(\Delta_0-k<0\), so it cannot
cover the trading cost. Against \(F_1\), every correct order earns at least \(\rho\kappa(\pi)\Delta_1-k>0\), so it
covers the cost for sure. Proposition~1(b) gives \(\Delta_2\ge\Delta_1\), so the right side also forces trade
against \(F_2\).

\emph{What fails without it.} Suppose \(\Delta_0>k\). Then the upper bound does not rule out informed orders
against \(F_0\), and the price can carry information there. Along the example's path the paper shows a strength
where the spread exceeds \(k\) and two equilibria exist (Proposition~3, Chapter~\ref{ch:coexistence}). Suppose
instead \(k\ge\rho\kappa(\pi)\Delta_1\). Then the lower bound does not force full orders against \(F_1\). The proof
cannot exclude candidates in which the investor abstains or mixes, and its uniqueness argument breaks.

\subsection{(H3), the exclusion condition}

\emph{What it compares.} (H3) reads \(B_2(M(\pi))<c_H\). It compares the high cost with gross profit at the top
posterior against the very strong incumbent.

\emph{Why the proof needs it.} Against \(F_2\) the investor trades fully, so every belief at a price lies between
\(m(\pi)\) and \(M(\pi)\). Gross profit never falls with the belief, so at every price it is at most
\(B_2(M(\pi))\), which is below \(c_H\). No price recruits a high-cost challenger, and entry returns to \(\rho\).

\emph{What fails without it.} A purchase could recruit a high-cost challenger against \(F_2\) as well, and entry
need not fall. The proof uses (H3) only against \(F_2\), so parts (i) and (ii) hold without it.

\begin{watchout}
The main text of the paper names the exclusion bound when it explains (H3). The exclusion bound of
Chapter~\ref{ch:beliefs} needs \(\beta\le1/3\), because it must bound the posterior after the zero order under every
strategy. Proposition~2 allows \(0<\beta<1/2\). The proof of part (iii) avoids the zero-order bound: it first shows
that orders are full, so the posterior after the zero order is exactly \(\pi\). Then it applies the logic of the
exclusion bound to three known posteriors. Section~\ref{sec:ccc-proof-iii} gives the details.
\end{watchout}

\subsection{Each inequality has one job}

Appendix A.3 of the paper lists where the proof uses each of the five inequalities. Table~\ref{tab:ccc-jobs}
repeats the list.

\begin{table}[ht]
\centering
\small
\begin{tabular}{@{}>{\raggedright\arraybackslash}p{3.4cm}>{\raggedright\arraybackslash}p{11.4cm}@{}}
\toprule
Inequality & Where the proof uses it\\
\midrule
\(B_0(\pi)<c_H\) & Entry \(\rho\) against \(F_0\). With Proposition~1(a), no recruitment after the zero order against
\(F_1\).\\
\(c_H<B_1(M(\pi))\) & Recruitment after a purchase against \(F_1\).\\
\(\Delta_0<k\) & No informed order pays against \(F_0\), under every candidate.\\
\(k<\rho\kappa(\pi)\Delta_1\) & Full orders against \(F_1\), and with Proposition~1(b) against \(F_2\), under every
candidate.\\
\(B_2(M(\pi))<c_H\) & No recruitment against \(F_2\).\\
\bottomrule
\end{tabular}
\caption{The job of each inequality in (H1) to (H3).}
\label{tab:ccc-jobs}
\end{table}

The paper also states that (H1) to (H3) hold together on an open set of parameters. It argues this only for
uniform incumbents with the example's class values and reserve. There each condition is a strict inequality between
continuous functions of \((r_0,r_1,r_2,c_H,k,\rho,\pi,\lambda,\beta)\), so small changes of these parameters keep all
three. The set is nonempty because the example satisfies all three. Do not read the claim as a statement about
every pair of classes or every incumbent family. Chapter~\ref{ch:example} shows a section of this set.

\section{When the trading-cost window is nonempty}\label{sec:ccc-window}

(H2) asks for a trading cost between two numbers. Such a cost exists only if the left number lies below the right
one. The paper states the exact condition: the window is nonempty exactly when \(\Delta_0/\Delta_1<\rho\kappa(\pi)\), that is,
when competition widens the target-payoff spread enough from \(F_0\) to \(F_1\).

\begin{derivation}[the window is nonempty exactly when \(\Delta_0/\Delta_1<\rho\kappa(\pi)\)]
\emph{Step 1: a fact about intervals.} For two real numbers \(a\) and \(c\), some \(k\) satisfies \(a<k<c\) exactly
when \(a<c\). If \(a<c\), the midpoint works, because \(a<(a+c)/2<c\). If \(a\ge c\), every \(k>a\) also satisfies
\(k>c\), so no \(k\) lies below \(c\) and above \(a\).

\emph{Step 2: apply it.} Take \(a=\Delta_0\) and \(c=\rho\kappa(\pi)\Delta_1\). Some trading cost satisfies (H2)
exactly when
\begin{equation}\label{eq:ccc-window}
\Delta_0<\rho\kappa(\pi)\Delta_1 .
\end{equation}

\emph{Step 3: the signs.} Under the class ranking \(\Delta_0\ge0\) (Section~\ref{sec:ccc-setup}), and
Proposition~1(b) gives \(\Delta_1\ge\Delta_0\). If \(\Delta_1=0\), then \(\Delta_0=0\), and \eqref{eq:ccc-window} reads
\(0<0\), which is false. So a nonempty window needs \(\Delta_1>0\).

\emph{Step 4: divide.} With \(\Delta_1>0\),
\begin{align*}
\Delta_0<\rho\kappa(\pi)\Delta_1
&\iff \frac{\Delta_0}{\Delta_1}<\rho\kappa(\pi) \why{divide by \(\Delta_1>0\); direction kept}.
\end{align*}
Every cost in the window is positive, because \(k>\Delta_0\ge0\).
\end{derivation}

\subsection{No unit of money}

The paper notes a scale property of the model. Multiply every value, the reserve and every cost by the same
positive number \(s\). Then prices and payoffs are multiplied by \(s\), while orders, beliefs, entry and ownership do
not change. Under this rescale, \(\Delta_0\) and \(\Delta_1\) are both multiplied by \(s\), so their ratio does not
change. The entry floor \(\rho\) is a probability, and \(\kappa(\pi)\) is built from posteriors. So the condition
\(\Delta_0/\Delta_1<\rho\kappa(\pi)\) involves the two distributions, the classes, \(\rho\), the order slot and
\(\pi\), and no unit of money.

In words, the condition asks by what factor competition widens the spread. The size of the spread in money does
not enter. The trading cost \(k\) is an amount of money, and (H2) then places it inside the window.

\subsection{Why a prior away from one half narrows the window}

The left end \(\Delta_0\) of the window does not depend on the prior. The right end \(\rho\kappa(\pi)\Delta_1\) does,
through \(\kappa(\pi)\). The derivation below shows that \(\kappa(\pi)\) is largest at \(\pi=1/2\).

\begin{derivation}[\(\kappa(\pi)\) is largest at the prior one half]
\emph{Step 1: \(1-M(\pi)=m(1-\pi)\).} Start from the top posterior:
\begin{align*}
1-M(\pi)&=1-\frac{\pi(1+u)}{\pi(1+u)+1-\pi} \why{definition of \(M\)}\\
&=\frac{\pi(1+u)+1-\pi-\pi(1+u)}{\pi(1+u)+1-\pi} \why{common denominator}\\
&=\frac{1-\pi}{\pi(1+u)+1-\pi} \why{\(\pi(1+u)\) cancels}.
\end{align*}
Next, evaluate the bottom posterior at \(1-\pi\):
\begin{align*}
m(1-\pi)&=\frac{1-\pi}{(1-\pi)+[1-(1-\pi)](1+u)} \why{replace \(\pi\) by \(1-\pi\) in \(m\)}\\
&=\frac{1-\pi}{(1-\pi)+\pi(1+u)} \why{\(1-(1-\pi)=\pi\)}.
\end{align*}
The two results have the same numerator and the same denominator, so \(1-M(\pi)=m(1-\pi)\).

\emph{Step 2: \(m\) rises strictly with the prior.} The complement of \(m(\pi)\) is
\begin{align*}
1-m(\pi)&=\frac{\pi+(1-\pi)(1+u)-\pi}{\pi+(1-\pi)(1+u)}=\frac{(1-\pi)(1+u)}{\pi+(1-\pi)(1+u)} \why{common denominator}.
\end{align*}
So the odds of \(m(\pi)\) are
\begin{align*}
\frac{m(\pi)}{1-m(\pi)}&=\frac{\pi}{(1-\pi)(1+u)} \why{the denominators cancel}\\
&=\frac{1}{1+u}\cdot\frac{\pi}{1-\pi} \why{rearrange}.
\end{align*}
The prior odds \(\pi/(1-\pi)\) rise strictly with \(\pi\). A probability rises strictly with its odds
(Section~\ref{sec:toolkit-odds}). So \(m(\pi)\) rises strictly with \(\pi\).

\emph{Step 3: \(\kappa\) as one bottom posterior.} By Step 1, \(\kappa(\pi)=\min\{m(\pi),m(1-\pi)\}\). Since \(m\)
rises, the smaller argument gives the smaller value:
\[
\kappa(\pi)=m\bigl(\min\{\pi,1-\pi\}\bigr).
\]

\emph{Step 4: the largest value.} The number \(\min\{\pi,1-\pi\}\) is at most \(1/2\), with equality only at
\(\pi=1/2\). Since \(m\) rises strictly, \(\kappa(\pi)\le m(1/2)\), with equality only at \(\pi=1/2\). The value is
\begin{align*}
m(1/2)&=\frac{1/2}{1/2+(1/2)(1+u)} \why{definition of \(m\)}\\
&=\frac{1}{1+(1+u)} \why{multiply above and below by \(2\)}\\
&=\frac{1}{2+u}.
\end{align*}
\end{derivation}

\emph{What this says.} The right end of the window is largest at the prior one half. It shrinks when the prior
moves away from one half in either direction, while the left end stays where it is. So a prior away from one half
narrows the window. It follows that, when \(\Delta_0>0\), a prior far enough from one half closes it, because
\(\kappa(\pi)\) goes to zero as \(\pi\) goes to zero or one.

\begin{intuition}[the favored class binds]
Suppose the prior is near one. The market already leans toward class \(H\). A purchase moves the belief to the top
posterior, which is then close to one. Class \(H\)'s return is entry times the spread times \(1-\mu(1)\), and that
last factor is small. The purchase tells the market little that it does not already believe. The paper puts it
this way: the price already leans toward the favored class.
\end{intuition}

Two more extreme cases help.
\begin{itemize}
\item \emph{The entry floor near zero.} As \(\rho\) falls toward zero, the right end \(\rho\kappa(\pi)\Delta_1\)
  falls toward zero, and the window shrinks; when \(\Delta_0>0\), it closes. The floor guarantees the investor a
  return under every candidate, because zero-cost challengers prepare after every order.
\item \emph{More informed presence.} The value \(\kappa(1/2)=1/(2+u)\) falls as the mass ratio \(u\) rises: a more
  informative order leaves less doubt behind. The guide adds that \(M(\pi)\), with odds \((1+u)\pi/(1-\pi)\), rises with
  \(u\), which helps (H1). So (H1) wants an informative purchase, and (H2) wants some doubt after it.
  Chapter~\ref{ch:example} shows the band of \(u\) where both hold in the example.
\end{itemize}

The guide adds one more consequence. Step 4 gives \(\kappa(\pi)\le1/(2+u)<1/2\), and \(\rho<1\), so
\(\rho\kappa(\pi)<1/2\). A nonempty window then needs \(\Delta_0/\Delta_1<1/2\), that is, \(\Delta_1>2\Delta_0\). When
the weak incumbent's spread is positive, the strong incumbent's spread must be more than twice as large.

\begin{checkyourself}
Let \(\pi>1/2\). Which of the two correct orders, the purchase of class \(H\) or the sale of class \(L\), sets
\(\kappa(\pi)\)? Explain in words why it is the order of the class that the prior favors.
\label{q:ccc-window}
\end{checkyourself}

\section{The proof}\label{sec:ccc-proof}

The proof in Appendix A.3 of the paper has seven steps. Step~0 collects inequalities that (H1) to (H3) imply. Step~1
solves the weak incumbent. Steps~2 to 4 solve the strong incumbent. Step~5 solves the very strong incumbent. Step~6
explains why the tie rules do not matter.

\begin{keyidea}
In each economy the proof fixes the investor's orders first, with a bound that holds for every candidate. The
orders give the posteriors. The posteriors give entry. Entry and the posteriors give the prices. Because the first
bound holds for every candidate, the tuple that comes out is the only equilibrium.
\end{keyidea}

At the end of each part the proof checks (E1) to (E6) at the tuple. Uniqueness gives at most one equilibrium, and
the check gives at least one.

\subsection{Step 0: inequalities that the conditions imply}\label{sec:ccc-proof-0}

Seven facts are used again and again below. Each one comes from a condition and an earlier result.

\begin{derivation}[Step 0]
\emph{(a) Gross profit at the prior is lower against \(F_1\).} Proposition~1(a) gives \(g_{L,1}\le g_{L,0}\) and
\(g_{H,1}\le g_{H,0}\). So
\begin{align*}
B_1(\pi)&=(1-\pi)\,g_{L,1}+\pi\,g_{H,1} \why{definition of \(B\)}\\
&\le(1-\pi)\,g_{L,0}+\pi\,g_{H,0} \why{Proposition~1(a); weights \(\ge0\)}\\
&=B_0(\pi).
\end{align*}
The same lines hold at every belief, so \(B_1\le B_0\) at every point. With (H1) this gives
\begin{equation}\label{eq:ccc-step0a}
B_1(\pi)\le B_0(\pi)<c_H<B_1(M).
\end{equation}

\emph{(b) The top posterior lies above the prior, and the bottom posterior below it.} The odds of \(M\) are
\((1+u)\,\pi/(1-\pi)\), and the odds of \(m\) are \(\pi/[(1-\pi)(1+u)]\). Since \(u=\lambda/b>0\), the first odds are
larger than the prior odds and the second are smaller. So \(m<\pi<M\).

\emph{(c) Class \(H\) earns more than class \(L\) against \(F_1\).}
\begin{align*}
0&<B_1(M)-B_1(\pi) \why{\eqref{eq:ccc-step0a}}\\
&=(M-\pi)\,(g_{H,1}-g_{L,1}) \why{\(B\) is linear in the belief}.
\end{align*}
Since \(M-\pi>0\) by (b), divide by \(M-\pi\) to get \(g_{H,1}-g_{L,1}>0\).

\emph{(d) The threshold against \(F_1\) lies between the prior and the top posterior.} By (c), the threshold
\(\tau_1=(c_H-g_{L,1})/(g_{H,1}-g_{L,1})\) exists, \(B_1\) rises strictly, and \(B_1(\tau_1)=c_H\). Then
\eqref{eq:ccc-step0a} reads \(B_1(\pi)<B_1(\tau_1)<B_1(M)\), and a function that rises strictly keeps the order of
its arguments. So \(\pi<\tau_1<M\).

\emph{(e) The spread is positive and large against \(F_1\) and \(F_2\).} Both \(m\) and \(1-M\) are positive, so
\(\kappa>0\). From the right side of (H2),
\begin{align*}
\Delta_1&>\frac{k}{\rho\kappa} \why{divide \(k<\rho\kappa\Delta_1\) by \(\rho\kappa>0\)}\\
&>0 \why{\(k>0\)}.
\end{align*}
\(F_2\) strengthens \(F_1\), so Proposition~1(b) gives \(\Delta_2\ge\Delta_1\). Hence
\begin{align*}
\rho\kappa\Delta_2&\ge\rho\kappa\Delta_1 \why{\(\Delta_2\ge\Delta_1\), \(\rho\kappa>0\)}\\
&>k \why{(H2)}.
\end{align*}

\emph{(f) Gross profit against \(F_2\) is below \(c_H\) at every belief up to \(M\).} Under the class ranking,
\(g_{H,2}\ge g_{L,2}\) (Section~\ref{sec:ccc-setup}). For every \(\mu\le M\),
\begin{align*}
B_2(\mu)&=B_2(M)-(M-\mu)(g_{H,2}-g_{L,2}) \why{\(B\) is linear in the belief}\\
&\le B_2(M) \why{\(M-\mu\ge0\), \(g_{H,2}-g_{L,2}\ge0\)}\\
&<c_H \why{(H3)}.
\end{align*}

\emph{(g) The entry floor.} A zero-cost challenger prepares after every price, so in every candidate of every
economy \(e(x)=\rho+(1-\rho)\zeta(P(x))\ge\rho\).
\end{derivation}

Notice what (f) does not use. It does not need the threshold \(\tau_2\), which may not exist, because the class
ranking gives only \(g_{H,2}\ge g_{L,2}\), with equality allowed.

\subsection{Part (i): the weak incumbent}\label{sec:ccc-proof-i}

Fix any candidate against \(F_0\). Nothing about it is assumed beyond (E1) to (E4) and (E6).

\begin{derivation}[Step 1: the orders, beliefs, entry and price against \(F_0\)]
\emph{1. No correct order pays.} The trading bound gives \(K_H\le\Delta_0\) and \(K_L\le\Delta_0\). For the purchase,
\begin{align*}
K_H&=e(1)\,\Delta_0\,[1-\mu(1)] \why{trading return}\\
&\le1\cdot\Delta_0\cdot1 \why{\(e(1)\le1\), \(0\le1-\mu(1)\le1\), \(\Delta_0\ge0\)}\\
&=\Delta_0 .
\end{align*}
So class \(H\)'s purchase earns
\begin{align*}
U_H(1)&=K_H-k \why{trading payoff}\\
&\le\Delta_0-k \why{\(K_H\le\Delta_0\)}\\
&<0=U_H(0) \why{left side of (H2)}.
\end{align*}
The same lines with \(K_L=e(-1)\Delta_0\mu(-1)\le\Delta_0\) give \(U_L(-1)<0=U_L(0)\).

\emph{2. No wrong-signed order pays.} For class \(H\)'s sale,
\begin{align*}
U_H(-1)&=-e(-1)\,\Delta_0\,[1-\mu(-1)]-k \why{trading payoff}\\
&\le-k<0 \why{the first term is \(\le0\)}.
\end{align*}
In the same way \(U_L(1)=-e(1)\Delta_0\mu(1)-k\le-k<0\).

\emph{3. Both classes abstain.} By 1 and 2, abstention is the unique best order of each class. (E5) then gives
\(\sigma_H=\sigma_L=\delta_0\). No investor is indifferent, so both investor tie rules give this answer.

\emph{4. Every posterior is the prior.} With \(\sigma_H=\sigma_L=\delta_0\), each order has the same probability in
both classes:
\[
f_H(x)=f_L(x)=(1-\lambda)\eta(x)+\lambda\1\{x=0\}.
\]
Bayes' rule in odds form gives
\begin{align*}
\frac{\mu(x)}{1-\mu(x)}&=\frac{\pi}{1-\pi}\cdot\frac{f_H(x)}{f_L(x)} \why{(E1), odds form}\\
&=\frac{\pi}{1-\pi} \why{\(f_H(x)=f_L(x)\)}.
\end{align*}
Equal odds mean equal probabilities, so \(\mu(x)=\pi\) after every order.

\emph{5. The belief at every price is the prior.} Take any price \(\bar P\), and let \(X(\bar P)\) be the set of
orders with this price. (E2) gives
\begin{align*}
\frac{\bar\mu(\bar P)}{1-\bar\mu(\bar P)}
&=\frac{\pi}{1-\pi}\cdot\frac{\sum_{x\in X(\bar P)}f_H(x)}{\sum_{x\in X(\bar P)}f_L(x)} \why{(E2), odds form}\\
&=\frac{\pi}{1-\pi} \why{the two sums are equal}.
\end{align*}
So \(\bar\mu(\bar P)=\pi\) at every price, whatever orders share it.

\emph{6. Only the zero-cost challenger prepares.} At every price the belief is \(\pi\), and \(B_0(\pi)<c_H\) by the left
side of (H1). (E3) gives \(\zeta(\bar P)=0\), with no tie. So \(e(x)=\rho+(1-\rho)\cdot0=\rho\) after every order.

\emph{7. One price for all orders.} The price formula with \(e(x)=\rho\) and \(\mu(x)=\pi\) gives
\[
P(x)=t_{0,0}+\rho\,[t_{L,0}-t_{0,0}+\Delta_0\pi]=:\bar P_0\qquad\text{for every order }x.
\]

\emph{8. Entry.} Entry is \(\rho\) after every order, so total entry is \(\rho\).

\emph{9. High-class ownership.} Preparation depends only on the order and the cost. Given the class, both are
independent of the values \((\theta,R)\). So
\begin{align*}
\Own_{H,0}&=\Pr(j=H)\cdot\Pr(\text{prepare}\mid H)\cdot\Pr(\theta>R\mid H,\text{prepare}) \why{chain rule}\\
&=\pi\cdot\rho\cdot\Pr(\theta>R\mid H) \why{entry \(\rho\); independence}\\
&=\pi\rho\,\E_H[F_0(\theta)] \why{\(\Pr(R<\theta\mid\theta)=F_0(\theta)\)}\\
&=\pi\rho\,W_H(F_0) \why{class ranking: \(\theta\ge p\)}.
\end{align*}
\end{derivation}

Steps 1 to 9 show that every equilibrium against \(F_0\) is the tuple below, with truthful bids. The appendix
calls it \(\mathcal E_{\rm N}(F_0)\).
\[
\sigma_H=\sigma_L=\delta_0,\qquad \mu(x)=\pi,\qquad P(x)=\bar P_0,\qquad \bar\mu(\bar P_0)=\pi,\qquad
\zeta(\bar P_0)=0.
\]

\emph{The tuple is an equilibrium.} (E1): \(f_H=f_L\), so every posterior is \(\pi\). (E2): the single price pools
all orders, and \(\sum_xf_H(x)=\sum_xf_L(x)\), so its belief is \(\pi\). (E3): \(B_0(\pi)<c_H\), strictly, so
\(\zeta=0\) is optimal. (E4): with \(e\equiv\rho\), the expected terminal payoff after every order is
\(t_{0,0}+\rho\,w_0(\pi)=\bar P_0\). (E5): \(K_H=\rho\Delta_0(1-\pi)\le\Delta_0<k\) and
\(K_L=\rho\Delta_0\pi\le\Delta_0<k\), so abstention beats both correct orders, and wrong-signed orders earn at most
\(-k\). (E6): truthful bids are weakly dominant (Chapter~\ref{ch:auction}).

\begin{intuition}[a silent price]
No informed order can pay against the weak incumbent, whatever the market makers believe and whatever entry
follows. The best case for a purchase gives the whole spread, and the spread is below the trading cost. With no
informed trade, every order is liquidity noise, and the belief stays at the prior. At the prior the high cost is
not covered.
\end{intuition}

An extreme case tests the argument. If \(\Delta_0=0\), the argument works unchanged: it never uses that the price
reveals the posterior, and with a zero spread the price would not move with the posterior anyway.

\begin{checkyourself}
Suppose that, against the weak incumbent, market makers expected class \(H\) to buy, and set a high price after a
purchase. Use the trading bound to explain why class \(H\) would still not buy. Why does the argument not depend on
what the market makers expect?
\label{q:ccc-silent}
\end{checkyourself}

\subsection{Part (ii), Step 2: the orders against the strong incumbent}\label{sec:ccc-proof-ii}

Fix any candidate against \(F_1\). The argument mirrors Step~1, with the lower end of the trading bound in place of
the upper end.

\begin{derivation}[Step 2: class \(H\) buys and class \(L\) sells]
\emph{1. The purchase covers the trading cost.}
\begin{align*}
K_H&=e(1)\,\Delta_1\,[1-\mu(1)] \why{trading return}\\
&\ge\rho\,\Delta_1\,[1-M] \why{\(e(1)\ge\rho\), \(\mu(1)\le M\), \(\Delta_1>0\)}\\
&\ge\rho\kappa\Delta_1 \why{\(1-M\ge\kappa\)}\\
&>k \why{right side of (H2)}.
\end{align*}
The bound \(\mu(1)\le M\) is Lemma~A.1(b). It holds under every investor strategy, because it concerns a nonzero
order.

\emph{2. The sale covers the trading cost.}
\begin{align*}
K_L&=e(-1)\,\Delta_1\,\mu(-1) \why{trading return}\\
&\ge\rho\,\Delta_1\,m \why{\(e(-1)\ge\rho\), \(\mu(-1)\ge m\), \(\Delta_1>0\)}\\
&\ge\rho\kappa\Delta_1 \why{\(m\ge\kappa\)}\\
&>k \why{right side of (H2)}.
\end{align*}

\emph{3. The best orders.} For class \(H\),
\[
U_H(1)=K_H-k>0=U_H(0)>-k\ge U_H(-1).
\]
So the purchase is class \(H\)'s unique best order. For class \(L\),
\[
U_L(-1)=K_L-k>0=U_L(0)>-k\ge U_L(1).
\]
So the sale is class \(L\)'s unique best order. (E5) gives \(\sigma_H=\delta_1\) and \(\sigma_L=\delta_{-1}\). No
investor is indifferent, so both investor tie rules give this answer.
\end{derivation}

These bounds use only the posteriors after the two nonzero orders, so Step 2 holds for every \(\beta<1/2\).

\subsection{Part (ii), Step 3: posteriors, entry and prices}

\begin{derivation}[Step 3: the posteriors]
With \(\sigma_H=\delta_1\) and \(\sigma_L=\delta_{-1}\), a sale comes only from a liquidity trader in class \(H\), and
from a liquidity trader or the investor in class \(L\). The same argument applies to the other orders:
\[
f_H=(b,\;b_0,\;b+\lambda),\qquad f_L=(b+\lambda,\;b_0,\;b)\qquad\text{on the orders }(-1,0,1).
\]
The likelihood ratios are
\begin{align*}
\frac{f_H(1)}{f_L(1)}&=\frac{b+\lambda}{b}=1+\frac{\lambda}{b}=1+u \why{definition of \(u\)},\\
\frac{f_H(0)}{f_L(0)}&=\frac{b_0}{b_0}=1,\\
\frac{f_H(-1)}{f_L(-1)}&=\frac{b}{b+\lambda}=\frac{1}{1+u} \why{divide above and below by \(b\)}.
\end{align*}
After a purchase, the odds are \((1+u)\,\pi/(1-\pi)\). A probability with odds \(o\) equals \(o/(1+o)\), so
\begin{align*}
\mu(1)&=\frac{(1+u)\pi/(1-\pi)}{1+(1+u)\pi/(1-\pi)} \why{probability from odds}\\
&=\frac{\pi(1+u)}{(1-\pi)+\pi(1+u)} \why{multiply above and below by \(1-\pi\)}\\
&=M .
\end{align*}
After no order the odds are the prior odds, so \(\mu(0)=\pi\). After a sale the odds are
\(\pi/[(1-\pi)(1+u)]\), and the same two lines give \(\mu(-1)=\pi/[\pi+(1-\pi)(1+u)]=m\). This is Lemma~A.1(d).
\end{derivation}

\emph{The challenger reads the posterior from the price.} The three posteriors \(m<\pi<M\) differ (Step~0(b)).
Suppose two orders \(x\) and \(x'\) had one price. Entry depends on the order only through the price, so
\(e(x)=e(x')=:\bar e\ge\rho>0\). The price formula then gives
\[
0=P(x)-P(x')=\bar e\,\Delta_1\,[\mu(x)-\mu(x')].
\]
Since \(\bar e>0\) and \(\Delta_1>0\), this forces \(\mu(x)=\mu(x')\), which is false. So the three orders have three
different prices. Each price is reached by one order, and (E2) gives \(\bar\mu(P(x))=\mu(x)\).

\begin{derivation}[Step 3: entry after each order]
At each price the belief is the posterior of its order. (E3) compares gross profit at that belief with \(c_H\).
\begin{itemize}
\item After a purchase the belief is \(M\), and \(B_1(M)>c_H\) by (H1). So \(\zeta=1\) and \(e(1)=1\).
\item After no order the belief is \(\pi\), and \(B_1(\pi)<c_H\) by \eqref{eq:ccc-step0a}. So \(\zeta=0\) and
  \(e(0)=\rho\).
\item After a sale the belief is \(m\). Since \(g_{H,1}>g_{L,1}\) and \(m<\pi\),
  \begin{align*}
  B_1(m)&=B_1(\pi)-(\pi-m)(g_{H,1}-g_{L,1}) \why{\(B\) is linear}\\
  &<B_1(\pi)<c_H \why{both factors positive; \eqref{eq:ccc-step0a}}.
  \end{align*}
  So \(\zeta=0\) and \(e(-1)=\rho\). The weak inequality \(B_1(m)\le B_1(\pi)\) would already be enough, because
  \(B_1(\pi)<c_H\) is strict; the paper states the strict one, and the reason is \(g_{H,1}>g_{L,1}\) from Step~0(c).
\end{itemize}
Every comparison is strict, so the challenger's tie rule is never used. In terms of the threshold: \(m<\pi<\tau_1<M\)
by Step~0(d), so only a purchase reaches \(\tau_1\).
\end{derivation}

Table~\ref{tab:ccc-strong} collects the order probabilities, the posteriors and entry against \(F_1\).

\begin{table}[ht]
\centering
\begin{tabular}{@{}lccc@{}}
\toprule
Order \(x\) & sale, \(-1\) & no order, \(0\) & purchase, \(1\)\\
\midrule
Probability in class \(H\), \(f_H(x)\) & \(b\) & \(b_0\) & \(b+\lambda\)\\
Probability in class \(L\), \(f_L(x)\) & \(b+\lambda\) & \(b_0\) & \(b\)\\
Posterior \(\mu(x)\) & \(m\) & \(\pi\) & \(M\)\\
Entry \(e(x)\) & \(\rho\) & \(\rho\) & \(1\)\\
\bottomrule
\end{tabular}
\caption{The equilibrium against the strong incumbent, order by order.}
\label{tab:ccc-strong}
\end{table}

\begin{derivation}[Step 3: prices rise strictly with the order]
\emph{An identity for price differences.} For any two orders \(x\) and \(x'\), the price formula gives
\begin{align*}
P(x)-P(x')&=e(x)\,w(\mu(x))-e(x')\,w(\mu(x')) \why{price formula}\\
&=e(x)\,w(\mu(x))-e(x)\,w(\mu(x'))\\
&\qquad+e(x)\,w(\mu(x'))-e(x')\,w(\mu(x')) \why{add and subtract}\\
&=e(x)\,[w(\mu(x))-w(\mu(x'))]+[e(x)-e(x')]\,w(\mu(x')) \why{collect terms}\\
&=e(x)\,\DT\,[\mu(x)-\mu(x')]+[e(x)-e(x')]\,w(\mu(x')) \why{definition of \(w\)}.
\end{align*}
In the first line the level \(t_0\) cancels. The second line adds and subtracts \(e(x)\,w(\mu(x'))\). The result is
the increment identity of Appendix A.2 of the paper.

\emph{From a sale to no order.} Entry is \(\rho\) after both orders, so the second term is zero:
\begin{align*}
P(0)-P(-1)&=\rho\,\Delta_1\,(\pi-m) \why{identity with \(e(0)=e(-1)=\rho\)}\\
&>0 \why{\(\rho,\Delta_1>0\) and \(\pi>m\)}.
\end{align*}

\emph{From no order to a purchase.} Entry rises from \(\rho\) to \(1\):
\begin{align*}
P(1)-P(0)&=1\cdot\Delta_1\,(M-\pi)+(1-\rho)\,w_1(\pi) \why{identity with \(e(1)=1\), \(e(0)=\rho\)}\\
&>0 \why{first term \(>0\); second \(\ge0\)}.
\end{align*}
The second term is nonnegative because \(w_1(\pi)=(1-\pi)w_{L,1}+\pi w_{H,1}\) and \(w_{L,1},w_{H,1}\ge0\). So
\(P(-1)<P(0)<P(1)\).
\end{derivation}

Steps 2 and 3 show that every equilibrium against \(F_1\) is the tuple below, with truthful bids. The appendix
calls it \(\mathcal E_{\rm F}(F_1)\), and the guide reads the letter F as ``full orders''.
\[
\begin{aligned}
&\sigma_H=\delta_1,\quad \sigma_L=\delta_{-1},\quad (\mu(-1),\mu(0),\mu(1))=(m,\pi,M),\quad
(e(-1),e(0),e(1))=(\rho,\rho,1),\\
&P(x)=t_{0,1}+e(x)\,w_1(\mu(x)),\quad \bar\mu(P(x))=\mu(x),\quad \zeta(P(x))=\1\{x=1\}.
\end{aligned}
\]

\emph{The tuple is an equilibrium.} (E1): the order probabilities above give \((m,\pi,M)\). (E2): entry rises with
the posterior and lies in \([\rho,1]\), and \(\Delta_1>0\), so by the identity above the three prices rise strictly;
each price has one order, so \(\bar\mu(P(x))=\mu(x)\). (E3): \(B_1(M)>c_H\), \(B_1(\pi)<c_H\) and \(B_1(m)<c_H\), all
strict. (E4): the price formula with the stated entry. (E5): \(K_H=1\cdot\Delta_1(1-M)\ge\rho\kappa\Delta_1>k\) and
\(K_L=\rho\Delta_1m\ge\rho\kappa\Delta_1>k\), so \(U_H(1)>0=U_H(0)>U_H(-1)\) and \(U_L(-1)>0=U_L(0)>U_L(1)\). (E6):
truthful bids are weakly dominant.

\begin{intuition}[only a purchase recruits]
The investor's purchase is what lifts the belief to the top posterior. A purchase is still not proof of class
\(H\), because a liquidity trader buys with probability \(b\) in either class. But it is good enough news: gross
profit at \(M\) covers the high cost. No order and a sale leave the belief at or below the prior, where it does
not. So the price speaks after every order, and only one of its messages changes what the challenger does.
\end{intuition}

\subsection{Part (ii), Step 4: entry, ownership and the takeover share}

\begin{derivation}[Step 4: entry in each class and total entry]
Entry in class \(j\) adds up, over the three orders, the probability of the order in class \(j\) times the
probability of preparation after it. Given the order \(x\), preparation depends only on the price \(P(x)\) and the
cost, and the cost is independent of the class. So the probability of preparation given the class \(j\) and the
order \(x\) is \(e(x)\), the same in both classes. The three liquidity and informed masses add up to one:
\begin{align*}
2b+b_0+\lambda&=2(1-\lambda)\beta+(1-\lambda)(1-2\beta)+\lambda \why{definitions of \(b,b_0\)}\\
&=(1-\lambda)(2\beta+1-2\beta)+\lambda=1 \why{factor \(1-\lambda\)}.
\end{align*}
So \(b+b_0=1-(b+\lambda)\) and \((b+\lambda)+b_0=1-b\). For class \(H\),
\begin{align*}
e_H&=f_H(-1)\,e(-1)+f_H(0)\,e(0)+f_H(1)\,e(1) \why{total probability; \(\Pr(\text{prepare}\mid H,x)=e(x)\)}\\
&=b\rho+b_0\rho+(b+\lambda)\cdot1 \why{Table~\ref{tab:ccc-strong}}\\
&=\rho\,[1-(b+\lambda)]+(b+\lambda) \why{\(b+b_0=1-(b+\lambda)\)}\\
&=\rho+(1-\rho)(b+\lambda) \why{rearrange}.
\end{align*}
For class \(L\),
\begin{align*}
e_L&=(b+\lambda)\rho+b_0\rho+b\cdot1 \why{Table~\ref{tab:ccc-strong}}\\
&=\rho\,(1-b)+b \why{\((b+\lambda)+b_0=1-b\)}\\
&=\rho+(1-\rho)\,b \why{rearrange}.
\end{align*}
Total entry averages the two with the prior weights:
\begin{align}
\Ent_1&=\pi e_H+(1-\pi)e_L \why{total probability over classes}\notag\\
&=\rho+(1-\rho)\,[\pi(b+\lambda)+(1-\pi)b] \why{substitute; \(\pi+(1-\pi)=1\)}\notag\\
&=\rho+(1-\rho)(b+\pi\lambda). \whyinline{\(\pi b+(1-\pi)b=b\)}\label{eq:ccc-entry}
\end{align}
Since \(b=(1-\lambda)\beta>0\), \((1-\rho)(b+\pi\lambda)>0\), and \(\Ent_1>\rho\). Also
\(e_H-e_L=(1-\rho)\lambda>0\).
\end{derivation}

\emph{What this says.} Total entry rises above the floor by \((1-\rho)\) times the probability of a purchase. A
purchase recruits a high-cost challenger, and the high cost appears with probability \(1-\rho\). Class \(L\)'s entry
also rises, through liquidity purchases: the price cannot tell a liquidity purchase from an informed one. The extra
entry favors class \(H\), which reaches a purchase with probability \(b+\lambda\) against \(b\) for class \(L\).

\emph{Extreme cases.} If \(\rho\) is near one, almost every challenger has zero cost, and entry is near one in every
economy, so recruitment adds little. If the slot has almost no liquidity purchases (\(\beta\) near zero), class \(L\)
almost never reaches a purchase, and \(e_L\) is near \(\rho\). This case strains (H2), though: \(b\) near zero makes
\(u=\lambda/b\) large and \(\kappa(\pi)\) small, so the trading cost must then be small.

\emph{Ownership.} Write \(\Own_{j,i}\) for the probability that the challenger is of class \(j\) and acquires the
target against \(F_i\), and \(\pi_H=\pi\), \(\pi_L=1-\pi\). Preparation depends on the order and the cost, and given
the class both are independent of \((\theta,R)\). A prepared class-\(j\) challenger acquires the target when
\(\theta>R\), which has probability \(W_j(F_i)\). As in Step~1,
\[
\Own_{j,1}=\pi_j\,e_j\,W_j(F_1),\qquad \Own_{j,0}=\pi_j\,\rho\,W_j(F_0).
\]

\begin{derivation}[Step 4: when high-class ownership rises]
\begin{align*}
\Own_{H,1}>\Own_{H,0}
&\iff \pi\,e_H\,W_H(F_1)>\pi\,\rho\,W_H(F_0) \why{the two formulas}\\
&\iff e_H\,W_H(F_1)>\rho\,W_H(F_0) \why{divide by \(\pi>0\)}\\
&\iff [\rho+(1-\rho)(b+\lambda)]\,W_H(F_1)>\rho\,W_H(F_0) \why{substitute \(e_H\)}.
\end{align*}
\end{derivation}

Two effects meet in this condition. Entry in class \(H\) rises from \(\rho\) to \(e_H\). But class \(H\)'s chance to
beat the incumbent can fall, from \(W_H(F_0)\) to \(W_H(F_1)\), because the incumbent is stronger.
Section~\ref{sec:ccc-ownership} gives a sufficient condition for the first effect to win.

\emph{The takeover share needs positive acquisitions.} Class \(H\)'s takeover share against \(F_i\) is
\(\Own_{H,i}/(\Own_{H,i}+\Own_{L,i})\). The denominator must be positive in both economies. The paper argues this in
three steps, and the guide fills in the middle one.

\begin{derivation}[Step 4: the challenger acquires the target with positive probability]
\emph{1. Some class earns a positive gross profit against \(F_1\).} By (H1), \(B_1(M)>c_H>0\). Since
\(B_1(M)=(1-M)g_{L,1}+Mg_{H,1}\) is an average of the two gross profits, at least one of them is positive. Call
that class \(j\).

\emph{2. A positive gross profit means a positive chance to win.} Recall from Chapter~\ref{ch:auction} that a
challenger with value \(\theta\ge p\) has the gross profit \(g(\theta)=\int_p^\theta F(y)\,dy\). Suppose
\(W_j(F)=\E_j[F(\theta)]=0\). Then \(F(\theta)=0\) with probability one under \(G_j\). Since \(F\) never falls,
\(F(y)=0\) for every \(y\le\theta\), so \(g(\theta)=0\) with probability one, and \(g_j=\E_j[g(\theta)]=0\). Read
backward: \(g_j>0\) implies \(W_j(F)>0\). So \(W_j(F_1)>0\).

\emph{3. The weaker incumbent is easier to beat.} \(F_1\le F_0\) at every point, so
\(W_j(F_1)=\E_j[F_1(\theta)]\le\E_j[F_0(\theta)]=W_j(F_0)\), for both classes.

So for the class \(j\) of point 1, \(\Own_{j,1}=\pi_je_jW_j(F_1)>0\) and \(\Own_{j,0}=\pi_j\rho W_j(F_0)>0\). Both
denominators are positive.
\end{derivation}

\begin{derivation}[Step 4: when class \(H\)'s takeover share rises]
\emph{A fact about shares.} Let \(x,y,x',y'\ge0\) with \(x+y>0\) and \(x'+y'>0\). Then
\begin{align*}
\frac{x}{x+y}>\frac{x'}{x'+y'}
&\iff x\,(x'+y')>x'\,(x+y) \why{multiply by \((x+y)(x'+y')>0\)}\\
&\iff xx'+xy'>x'x+x'y \why{expand}\\
&\iff xy'>x'y \why{cancel \(xx'\)}.
\end{align*}
\emph{Apply it.} Take \(x=\Own_{H,1}\), \(y=\Own_{L,1}\), \(x'=\Own_{H,0}\) and \(y'=\Own_{L,0}\). Then
\begin{align*}
xy'&=\pi e_HW_H(F_1)\cdot(1-\pi)\rho W_L(F_0) \why{ownership formulas},\\
x'y&=\pi\rho W_H(F_0)\cdot(1-\pi)e_LW_L(F_1) \why{ownership formulas}.
\end{align*}
Divide both by \(\pi(1-\pi)\rho>0\). The share rises from \(F_0\) to \(F_1\) exactly when
\[
e_H\,W_H(F_1)\,W_L(F_0)>e_L\,W_H(F_0)\,W_L(F_1),
\]
which is \eqref{eq:ccc-share}.
\end{derivation}

\begin{derivation}[Step 4: the case where class \(H\) outbids the strong incumbent]
Let \(W_H(F_1)=1\).

\emph{1. Class \(H\) also outbids \(F_0\).} A probability is at most one, and point 3 above gives
\(W_H(F_0)\ge W_H(F_1)\). So \(1\ge W_H(F_0)\ge W_H(F_1)=1\), and \(W_H(F_0)=1\).

\emph{2. Ownership rises.} \(\Own_{H,1}=\pi e_H\cdot1>\pi\rho\cdot1=\Own_{H,0}\), because \(e_H-\rho=(1-\rho)(b+\lambda)>0\).

\emph{3. The share condition.} With \(W_H(F_1)=W_H(F_0)=1\), \eqref{eq:ccc-share} reads
\(e_H\,W_L(F_0)>e_L\,W_L(F_1)\).
\begin{itemize}
\item If \(W_L(F_0)>0\):
  \begin{align*}
  e_H\,W_L(F_0)&>e_L\,W_L(F_0) \why{\(e_H>e_L\), \(W_L(F_0)>0\)}\\
  &\ge e_L\,W_L(F_1) \why{\(W_L(F_0)\ge W_L(F_1)\), \(e_L>0\)}.
  \end{align*}
  So the share rises strictly.
\item If \(W_L(F_0)=0\): then \(0\le W_L(F_1)\le W_L(F_0)=0\), so \(\Own_{L,0}=\Own_{L,1}=0\). The share is
  \(\Own_{H,i}/(\Own_{H,i}+0)=1\) against both incumbents.
\end{itemize}
So the share never falls, and it rises strictly exactly when \(W_L(F_0)>0\).
\end{derivation}

\begin{intuition}[when class \(H\) always wins]
If a prepared class-\(H\) challenger beats the strong incumbent for sure, it beats the weak one too. Then the
stronger incumbent takes wins only from class \(L\). Class \(H\)'s acquisitions rise with its entry. Class \(L\)
gains entry too, but it also loses more auctions. So the takeovers that go to the challenger shift toward class
\(H\), as long as a low-class challenger can beat the weak incumbent at all.
\end{intuition}

\begin{watchout}
The clean rise of high-class ownership needs class \(H\) to outbid the stronger incumbent, \(W_H(F_1)=1\). Without it,
ownership rises only under the condition \(e_HW_H(F_1)>\rho W_H(F_0)\), and a stronger incumbent can lower class
\(H\)'s chance to win. Do not state ``high-class ownership rises'' without one of these two conditions.
\end{watchout}

\subsection{Part (iii): the very strong incumbent}\label{sec:ccc-proof-iii}

Fix any candidate against \(F_2\). The order of the argument matters here: orders first, then posteriors, then
beliefs at each price, then entry, then prices.

\begin{derivation}[Step 5: the very strong incumbent]
\emph{1. Full orders.} Step~0(e) gives \(\rho\kappa\Delta_2>k\). The two derivations of Step~2, with \(\Delta_2\) in
place of \(\Delta_1\), give
\begin{align*}
K_H&\ge\rho\,[1-M]\,\Delta_2\ge\rho\kappa\Delta_2>k \why{trading bound; Step~0(e)},\\
K_L&\ge\rho\,m\,\Delta_2\ge\rho\kappa\Delta_2>k \why{trading bound; Step~0(e)}.
\end{align*}
So \(\sigma_H=\delta_1\) and \(\sigma_L=\delta_{-1}\), under both investor tie rules.

\emph{2. The posteriors.} The order probabilities do not depend on the incumbent. So Step~3 applies unchanged, and
the posteriors after \((-1,0,1)\) are \((m,\pi,M)\).

\emph{3. Beliefs at each price.} The belief at a price averages the posteriors of the orders with this price. So it
lies between the smallest and the largest of them, and in particular in \([m,M]\).

\emph{4. No high-cost preparation.} Step~0(f) gives \(B_2(\mu)<c_H\) for every \(\mu\le M\). So gross profit at every
price belief is below \(c_H\). (E3) gives \(\zeta\equiv0\), with no tie, and \(e(x)=\rho\) after every order.

\emph{5. The three prices differ.} Entry is \(\rho\) after every order, so the increment identity of Step~3 loses its
second term:
\begin{align*}
P(x)-P(x')&=\rho\,\Delta_2\,[\mu(x)-\mu(x')] \why{identity with \(e(x)=e(x')=\rho\)}.
\end{align*}
Here \(\rho>0\), \(\Delta_2\ge\Delta_1>0\), and the three posteriors differ. So the three prices differ, each price
has one order, and \(\bar\mu(P(x))=\mu(x)\).

\emph{6. High-class ownership.} Entry is \(\rho\) after every order, so as in Step~1,
\(\Own_{H,2}=\pi\rho\,W_H(F_2)\).
\end{derivation}

So every equilibrium against \(F_2\) is the tuple below, with truthful bids. The appendix calls it
\(\mathcal E_{\rm F}(F_2)\).
\[
\begin{aligned}
&\sigma_H=\delta_1,\quad \sigma_L=\delta_{-1},\quad (\mu(-1),\mu(0),\mu(1))=(m,\pi,M),\\
&P(x)=t_{0,2}+\rho\,w_2(\mu(x)),\quad \bar\mu(P(x))=\mu(x),\quad \zeta\equiv0.
\end{aligned}
\]

\emph{The tuple is an equilibrium.} (E1): the order probabilities give \((m,\pi,M)\). (E2): \(\Delta_2>0\) and
\(m<\pi<M\), so the three prices differ, and each price has one order. (E3): \(B_2(\mu(x))<c_H\) at each posterior.
(E4): the price formula with \(e\equiv\rho\). (E5): \(K_H=\rho\Delta_2(1-M)\ge\rho\kappa\Delta_2>k\) and
\(K_L=\rho\Delta_2m\ge\rho\kappa\Delta_2>k\). (E6): truthful bids are weakly dominant.

\begin{intuition}[a price that speaks, with nothing good enough to say]
Against the very strong incumbent the investor trades, so the price moves with the order. The best news a price
can carry is a purchase, which lifts the belief to \(M\). Even at \(M\), gross profit is below the high cost. So the
price is informative, and entry stays at its floor.
\end{intuition}

The three prices differ even though entry is the same after every order. The reason is in the price formula:
\(P(x)=t_{0,2}+\rho\,[t_{L,2}-t_{0,2}+\Delta_2\mu(x)]\) moves with the posterior through the term
\(\rho\Delta_2\mu(x)\). This term is active because the zero-cost challenger prepares with probability \(\rho>0\) and
the spread is positive. With a zero spread the three prices would coincide, even with informed trade.

\begin{checkyourself}
The exclusion bound of Chapter~\ref{ch:beliefs} needs \(\beta\le1/3\), but Proposition~2 allows any \(\beta<1/2\).
Explain why the proof of part (iii) still works when \(1/3<\beta<1/2\). Which step would fail if you skipped point 1
of Step~5?
\label{q:ccc-beta}
\end{checkyourself}

\subsection{The last paragraph of the statement}

The proposition ends with four claims for each economy. The prices are the price formula at the posterior and
entry of each tuple. The price belief is the posterior: against \(F_0\) every posterior and the single price belief
equal \(\pi\), and against \(F_1\) and \(F_2\) each price has one order. Bids are truthful by (E6). For high-cost
preparation, solve \(e=\rho+(1-\rho)\zeta\) for \(\zeta\):
\begin{align*}
e-\rho&=(1-\rho)\,\zeta \why{subtract \(\rho\)}\\
\zeta&=\frac{e-\rho}{1-\rho} \why{divide by \(1-\rho>0\)}.
\end{align*}
So \(\zeta=0\) after every order against \(F_0\) and \(F_2\). Against \(F_1\), \(\zeta=1\) after a purchase and
\(\zeta=0\) otherwise.

Table~\ref{tab:ccc-tuples} puts the three tuples side by side.

\begin{table}[ht]
\centering
\small
\begin{tabular}{@{}lccc@{}}
\toprule
 & weak, \(F_0\) & strong, \(F_1\) & very strong, \(F_2\)\\
\midrule
Informed orders of classes \(H\) and \(L\) & \((0,0)\) & \((1,-1)\) & \((1,-1)\)\\
Posteriors after \((-1,0,1)\) & \((\pi,\pi,\pi)\) & \((m,\pi,M)\) & \((m,\pi,M)\)\\
Entry after \((-1,0,1)\) & \((\rho,\rho,\rho)\) & \((\rho,\rho,1)\) & \((\rho,\rho,\rho)\)\\
High-cost preparation & none & after a purchase & none\\
Prices & one price & \(P(-1)<P(0)<P(1)\) & \(P(-1)<P(0)<P(1)\)\\
Total entry & \(\rho\) & \(\rho+(1-\rho)(b+\pi\lambda)\) & \(\rho\)\\
High-class ownership & \(\pi\rho W_H(F_0)\) & \(\pi e_HW_H(F_1)\) & \(\pi\rho W_H(F_2)\)\\
\bottomrule
\end{tabular}
\caption{The unique equilibrium in each economy under (H1) to (H3). The prices against \(F_2\) rise with the order
because they rise with the posterior.}
\label{tab:ccc-tuples}
\end{table}

\subsection{Step 6: uniqueness under all four tie rules}\label{sec:ccc-proof-ties}

The four tie rules are the four pairs of an investor rule (may mix, or abstains) and a challenger rule (prepares,
or may prepare with any probability), as in Section~\ref{sec:ccc-setup}. A tie rule only says what a player does
when two choices give the same payoff. The investor rule matters only when
a correct order earns exactly the trading cost, \(K_j=k\) (Lemma~A.3(b)). The challenger rule matters only when gross
profit at the price belief equals the high cost, \(B(\bar\mu)=c_H\).

\begin{derivation}[Step 6: no player is ever indifferent]
\emph{The investor.} The comparisons with \(k\) hold for every candidate: \(K_j\le\Delta_0<k\) against \(F_0\)
(Step~1), \(K_j\ge\rho\kappa\Delta_1>k\) against \(F_1\) (Step~2), and \(K_j\ge\rho\kappa\Delta_2>k\) against \(F_2\)
(Step~5). Each is strict, so no candidate has an indifferent investor.

\emph{The high-cost challenger.} The investor's comparisons fix the orders in every equilibrium, and the orders
fix the beliefs. Against \(F_0\) the belief is \(\pi\), and \(B_0(\pi)<c_H\). Against \(F_1\) the beliefs are \(M\),
\(\pi\) and \(m\), with \(B_1(M)>c_H\), \(B_1(\pi)<c_H\) and \(B_1(m)<c_H\). Against \(F_2\) every belief lies in
\([m,M]\), where \(B_2<c_H\). Each comparison is strict, so no equilibrium has an indifferent high-cost challenger.
A candidate with other orders can have other beliefs, but it fails (E5), so its beliefs do not matter.
\end{derivation}

No tie rule is ever used, so all four pairs give the same unique outcome in each economy. The strict inequalities
in (H1) to (H3) do this work: they keep every player away from indifference. Propositions~3 and~4, by contrast,
use the default tie rules (Chapters~\ref{ch:coexistence} and~\ref{ch:research}).
Chapter~\ref{ch:example} shows an edge of the example where a purchase leaves the high-cost challenger exactly
indifferent, and the challenger's tie rule then matters.

\section{When high-class ownership rises}\label{sec:ccc-ownership}

High-class ownership multiplies entry in class \(H\), which rises against \(F_1\), by class \(H\)'s chance to beat
the incumbent, which can fall. The paper gives a sufficient condition for a rise, in two parts:
\begin{equation}\label{eq:ccc-sufficient}
(1+\rho)\,W_H(F_1)\ge2\rho\,W_H(F_0)>0\qquad\text{and}\qquad b+\lambda>\tfrac12 .
\end{equation}
In words, the stronger incumbent leaves the high class a good chance to win, and the high class places a purchase
more often than not.

\begin{derivation}[the sufficient condition implies \(e_HW_H(F_1)>\rho W_H(F_0)\)]
\emph{Step 1: entry in class \(H\) exceeds \((1+\rho)/2\).}
\begin{align*}
e_H&=\rho+(1-\rho)(b+\lambda) \why{Step~4}\\
&>\rho+(1-\rho)\cdot\tfrac12 \why{\(b+\lambda>\tfrac12\), \(1-\rho>0\)}\\
&=\frac{2\rho+1-\rho}{2}=\frac{1+\rho}{2} \why{common denominator}.
\end{align*}

\emph{Step 2: class \(H\) can beat the strong incumbent.} The first part of \eqref{eq:ccc-sufficient} gives
\((1+\rho)W_H(F_1)\ge2\rho W_H(F_0)>0\). Since \(1+\rho>0\), \(W_H(F_1)>0\).

\emph{Step 3: combine.}
\begin{align*}
e_H\,W_H(F_1)&>\frac{1+\rho}{2}\,W_H(F_1) \why{Step~1 times \(W_H(F_1)>0\)}\\
&\ge\frac{2\rho\,W_H(F_0)}{2} \why{first part of \eqref{eq:ccc-sufficient}, halved}\\
&=\rho\,W_H(F_0).
\end{align*}
This is the exact condition of Step~4, so high-class ownership rises.
\end{derivation}

\emph{What this says.} Divide the first part by \((1+\rho)W_H(F_0)>0\): it reads
\(W_H(F_1)/W_H(F_0)\ge2\rho/(1+\rho)\). Since \(\rho<1\), the right side is below one. So the strong incumbent may
lower class \(H\)'s chance to win, but not below the fraction \(2\rho/(1+\rho)\) of its chance against the weak
incumbent.

\emph{Extreme cases.} When \(\rho\) is near zero, the fraction is near zero: entry against the weak incumbent is
tiny, so almost any positive chance to win against \(F_1\) is enough. When \(\rho\) is near one, the fraction is near
one, so the chance to win must barely fall.

When class \(H\) outbids \(F_1\), its chance to win is one against both incumbents, and ownership rises with entry in
class \(H\) whatever the value of \(b+\lambda\) (Step~4). The stronger incumbent then takes wins only from the low
class. When a low-class challenger can beat the weak incumbent, the takeovers that go to the challenger shift toward
class \(H\). Section~\ref{sec:ccc-decisive} gives a condition on values and costs alone that delivers this case.

\begin{checkyourself}
Suppose \(b+\lambda\le1/2\), and keep the first part of \eqref{eq:ccc-sufficient}. Which step of the derivation
fails? Show that at the edge \((1+\rho)W_H(F_1)=2\rho W_H(F_0)\) the exact condition \(e_HW_H(F_1)>\rho W_H(F_0)\) then
fails.
\label{q:ccc-sufficient}
\end{checkyourself}

\section{Decisive advantage}\label{sec:ccc-decisive}

The cleanest ownership statements of part (ii) need \(W_H(F_1)=1\). The paper gives a condition on values and costs
alone that delivers it. A high-cost challenger prepares only when gross profit at its belief covers \(c_H\). When
\(g_H>g_L\), gross profit at any belief below one is less than class \(H\)'s gross profit \(g_H\). So recruitment
needs \(g_H>c_H\), and if class \(H\) earns more than \(c_H\) only against incumbents that it always beats,
recruitment implies that class \(H\) outbids the incumbent.

\begin{paperresult}{Lemma A.6 (Decisive advantage)}
Fix ranked finite-mean classes, a reserve \(p>0\), a high cost \(c_H>0\), and a family of continuous incumbent laws.
The family has \emph{decisive advantage} if, for every law \(F\) in it,
\[
g_H(F)>c_H\quad\Longrightarrow\quad W_H(F)=1.
\]
At every law in this family with \(\DT>0\) and \(g_H>g_L\), any equilibrium that recruits a high-cost challenger after
some order has \(W_H(F)=1\). This holds at every interior prior, under all four tie rules, both at fixed presence and
with research at a cap below one. Under (H1)--(H3), decisive advantage at \(F_1\) gives the ownership and
takeover-share conclusions under outbidding in Proposition~2(ii).

For point classes \(h>\ell\ge p\) and the full family \(F_r=U[0,r]\), \(r>0\), decisive advantage holds exactly when
\[
c_H\ge\frac{h^2-p^2}{2h}.
\]
\end{paperresult}

In words, a family has decisive advantage when class \(H\)'s gross profit exceeds the high cost only against
incumbents that class \(H\) beats for sure. Inside such a family, at a law with \(\DT>0\) and \(g_H>g_L\), wherever a
price recruits a high-cost challenger, class \(H\) outbids the incumbent, whatever the prior and the tie rule. Research is the topic of
Chapter~\ref{ch:research}; here it enters only through a cap below one on the investor's presence. The last part
says when the uniform family, with one value in each class, has decisive advantage.

\begin{derivation}[proof of the consequence]
\emph{Step 1: every belief is interior.} Every order has positive probability in both classes, because a liquidity
trader can place it: \(f_j(x)\ge(1-\lambda)\eta(x)>0\). With \(0<\pi<1\), Bayes' rule gives \(0<\mu(x)<1\) after every
order. A price belief averages such posteriors, so it also lies strictly between zero and one. With research the
same holds, because presence stays below a cap below one, so liquidity orders keep every order possible in both
classes.

\emph{Step 2: recruitment needs gross profit at least \(c_H\).} Suppose a high-cost challenger prepares with positive
probability after some order, at a price with belief \(\bar\mu\). If \(B(\bar\mu)<c_H\), (E3) would give \(\zeta=0\). So
\(B(\bar\mu)\ge c_H\), under either challenger tie rule.

\emph{Step 3: the gross profit of class \(H\) exceeds \(c_H\).}
\begin{align*}
c_H&\le B(\bar\mu) \why{Step~2}\\
&=g_L+\bar\mu\,(g_H-g_L) \why{definition of \(B\)}\\
&=g_H-(1-\bar\mu)(g_H-g_L) \why{rearrange}\\
&<g_H \why{\(1-\bar\mu>0\), \(g_H-g_L>0\)}.
\end{align*}

\emph{Step 4: decisive advantage.} Step~3 gives \(g_H(F)>c_H\). Decisive advantage then gives \(W_H(F)=1\).
\end{derivation}

The strict inequality in Step~3 turns the weak condition \(c_H\le B(\bar\mu)\) into the strict condition \(g_H>c_H\)
that the definition needs. The four steps never use \(\DT>0\). The guide keeps this hypothesis in the statement
because the paper states it.

\begin{derivation}[why decisive advantage at \(F_1\) gives \(W_H(F_1)=1\) under (H1)]
Step~0 gives \(\Delta_1>0\) and \(g_{H,1}>g_{L,1}\), and part (ii) shows that a purchase recruits a high-cost challenger
against \(F_1\). So the consequence applies at \(F_1\). You can also see it directly from (H1):
\begin{align*}
c_H&<B_1(M) \why{right side of (H1)}\\
&=g_{H,1}-(1-M)(g_{H,1}-g_{L,1}) \why{rearrange \(B_1\)}\\
&<g_{H,1} \why{\(M<1\), \(g_{H,1}>g_{L,1}\)}.
\end{align*}
If \(F_1\) belongs to a family with decisive advantage, \(g_{H,1}>c_H\) gives \(W_H(F_1)=1\). Step~4 then applies: class
\(H\) also outbids \(F_0\), high-class ownership rises from \(\pi\rho\) to \(\pi e_H\), and class \(H\)'s takeover share
does not fall, and rises strictly exactly when \(W_L(F_0)>0\).
\end{derivation}

\begin{derivation}[the uniform family]
Let class \(H\) have the single value \(h\), with \(h>\ell\ge p\), and let \(F_r\) be uniform on \([0,r]\):
\(F_r(y)=y/r\) for \(0\le y\le r\), and \(F_r(y)=1\) for \(y\ge r\).

\emph{The chance to win.} Since \(h\ge p\), \(W_H(F_r)=F_r(h)\). So \(W_H(F_r)=1\) when \(r\le h\), and
\(W_H(F_r)=h/r<1\) when \(r>h\).

\emph{The gross profit for \(r>h\).} With a single value \(h\), \(1-G_H(y)=1\) for \(y<h\) and \(0\) for \(y\ge h\). So
\begin{align*}
g_H(F_r)&=\int_p^\infty[1-G_H(y)]\,F_r(y)\,dy \why{\eqref{eq:ccc-prop1}}\\
&=\int_p^h F_r(y)\,dy \why{\(1-G_H=\1\{y<h\}\)}\\
&=\int_p^h\frac yr\,dy \why{\(y\le h<r\)}\\
&=\frac1r\cdot\frac{h^2-p^2}{2} \why{\(\int y\,dy=y^2/2\)}\\
&=\frac{h^2-p^2}{2r}.
\end{align*}
This falls strictly as \(r\) rises, because \(h>p\). For every \(r>h\) it lies below \((h^2-p^2)/(2h)\), its value at
\(r=h\), and it comes as close as you like to that value as \(r\) falls toward \(h\).

\emph{If \(c_H\ge(h^2-p^2)/(2h)\), decisive advantage holds.} For \(r>h\),
\(g_H(F_r)<(h^2-p^2)/(2h)\le c_H\). So \(g_H(F_r)>c_H\) can hold only when \(r\le h\), and there \(W_H(F_r)=1\).

\emph{If \(c_H<(h^2-p^2)/(2h)\), decisive advantage fails.} For \(r>h\),
\begin{align*}
g_H(F_r)>c_H&\iff\frac{h^2-p^2}{2r}>c_H \why{formula for \(r>h\)}\\
&\iff r<\frac{h^2-p^2}{2c_H} \why{multiply by \(r/c_H>0\)}.
\end{align*}
The bound \((h^2-p^2)/(2c_H)\) exceeds \(h\), because \(c_H<(h^2-p^2)/(2h)\). So every \(r\) strictly between \(h\) and
\((h^2-p^2)/(2c_H)\) has \(g_H(F_r)>c_H\) and \(W_H(F_r)=h/r<1\). This law breaks the definition.

The two cases together prove the equivalence.
\end{derivation}

\emph{What this says.} A high cost screens the incumbents against which a high-cost challenger can be recruited.
When the high cost is at least \((h^2-p^2)/(2h)\), only incumbents weak enough for class \(H\) to beat for sure leave
class \(H\) a gross profit above it. A low high cost lets recruitment happen against incumbents that sometimes beat
class \(H\). Chapter~\ref{ch:example} uses this condition in the example.

\begin{checkyourself}
The proof of the consequence uses a belief strictly below one. Suppose, against the model, that some price could
reveal class \(H\) for sure. Explain which step would break, and why decisive advantage might then say nothing.
\label{q:ccc-decisive}
\end{checkyourself}

\section{What the result says and what it does not say}\label{sec:ccc-reading}

\subsection{What it says}

Across the three economies, entry first rises and then falls as the incumbent strengthens. The rise comes from
information: the strong incumbent widens the spread enough for informed trade to pay, and a purchase then lifts the
belief high enough to recruit a high-cost challenger. The fall comes from profit: against the very strong incumbent,
no belief that a price can produce covers the high cost. The market makers anticipate the preparation that each
price induces, so the reversal holds under competitive prices.

The paper's introduction adds a sharper point: entry rises even though acquisition profit falls at the prior and at
every other fixed belief that gives the high class positive weight. Proposition~1(a) alone gives only a weak fall.
The guide derives the strict fall from (H2).

\begin{derivation}[gross profit falls strictly at every belief \(\mu>0\) from \(F_0\) to \(F_1\)]
\emph{Step 1: the spread rises strictly.} Since \(\rho<1\) and \(\kappa<1/2\), \(\rho\kappa<1\). With \(\Delta_1>0\),
\begin{align*}
\Delta_0&<k \why{left side of (H2)}\\
&<\rho\kappa\Delta_1 \why{right side of (H2)}\\
&<\Delta_1 \why{\(\rho\kappa<1\), \(\Delta_1>0\)}.
\end{align*}

\emph{Step 2: the spread change as an integral.} Under the class ranking the reserve term in \eqref{eq:ccc-prop1}
is zero, so
\begin{align*}
\Delta_1-\Delta_0&=\int_p^\infty[1-F_1(y)]\,D(y)\,dy-\int_p^\infty[1-F_0(y)]\,D(y)\,dy \why{\eqref{eq:ccc-prop1}}\\
&=\int_p^\infty[F_0(y)-F_1(y)]\,D(y)\,dy \why{one integral}.
\end{align*}
By Step~1 this integral is positive.

\emph{Step 3: compare the weights.} Since \(G_L(y)\le1\), \(D(y)=G_L(y)-G_H(y)\le1-G_H(y)\). Also
\(F_0(y)-F_1(y)\ge0\). So
\begin{align*}
\int_p^\infty[F_0(y)-F_1(y)]\,[1-G_H(y)]\,dy&\ge\int_p^\infty[F_0(y)-F_1(y)]\,D(y)\,dy \why{larger weight}\\
&>0 \why{Step~2}.
\end{align*}

\emph{Step 4: class \(H\)'s gross profit falls strictly.}
\begin{align*}
g_{H,0}-g_{H,1}&=\int_p^\infty[1-G_H(y)]\,F_0(y)\,dy-\int_p^\infty[1-G_H(y)]\,F_1(y)\,dy \why{\eqref{eq:ccc-prop1}}\\
&=\int_p^\infty[1-G_H(y)]\,[F_0(y)-F_1(y)]\,dy>0 \why{Step~3}.
\end{align*}

\emph{Step 5: gross profit at every belief \(\mu>0\).}
\begin{align*}
B_0(\mu)-B_1(\mu)&=(1-\mu)(g_{L,0}-g_{L,1})+\mu\,(g_{H,0}-g_{H,1}) \why{definition of \(B\)}\\
&>0 \why{first term \(\ge0\); second \(>0\)}.
\end{align*}
\end{derivation}

So deterrence is at work at every belief that gives class \(H\) positive weight, and entry still rises. If the
challenger's belief were held fixed, the stronger incumbent could not raise entry. The rise needs the change in
information.

\subsection{What it does not say}

\begin{itemize}
\item \emph{It is a comparison across economies.} Each economy is solved again, with its own prices, beliefs and
  entry. Inside one economy, a deviation by the investor faces that economy's fixed prices and entry. The result does
  not describe how one market reacts over time when its incumbent becomes stronger.
\item \emph{It says nothing about the strengths between the three economies.} Between the weak and the strong
  incumbent, a quiet market and an informative market can both be equilibria. Chapter~\ref{ch:coexistence} describes
  where they coexist.
\item \emph{It does not say that entry rises with every strengthening.} The rise from \(F_0\) to \(F_1\) needs (H1)
  and (H2). The fall from \(F_1\) to \(F_2\) needs (H3) as well. Without (H3), parts (i) and (ii) still hold, but
  entry need not fall again.
\item \emph{It does not say that the price reveals the class.} In each of the three equilibria every posterior lies
  between \(m(\pi)\) and \(M(\pi)\), because a liquidity trader can place every order.
\item \emph{It does not say that high-class ownership always rises.} Ownership rises exactly when
  \(e_HW_H(F_1)>\rho W_H(F_0)\), and the clean rise needs \(W_H(F_1)=1\).
\item \emph{It does not say that a stronger incumbent raises entry at fixed information.} When the investor's orders
  stay fixed at their informative level, a stronger incumbent cannot raise entry. Chapter~\ref{ch:example} states
  this fixed-information control with its hypotheses.
\end{itemize}

In Chapter~\ref{ch:example}, Corollary~1 gives the slots and priors at which (H1) to (H3) hold together in the
example, with its entry and ownership values.

\begin{chaptersummary}
\begin{itemize}
\item Proposition~2 compares three economies that differ only in incumbent strength, and each has a unique
  equilibrium under all four tie rules.
\item Against the weak incumbent no informed order covers the trading cost, so the price is silent and only the
  zero-cost challenger prepares.
\item Against the strong incumbent a purchase lifts the belief to \(M(\pi)\) and recruits a high-cost challenger,
  so entry rises to \(\rho+(1-\rho)(b+\pi\lambda)\).
\item Against the very strong incumbent the price still speaks, but no belief at a price covers the high cost, so
  entry returns to \(\rho\).
\item (H1) sets entry through profit, (H2) sets trade through the two ends of the trading bound, and (H3) stops
  recruitment.
\item The window in (H2) is nonempty exactly when \(\Delta_0/\Delta_1<\rho\kappa(\pi)\), a condition with no unit of
  money.
\item High-class ownership rises exactly when \(e_HW_H(F_1)>\rho W_H(F_0)\), and under decisive advantage (H1) gives
  \(W_H(F_1)=1\) and the clean rise.
\end{itemize}
\end{chaptersummary}

\answersection

\answer{\ref{q:ccc-missing}}{Step~0(a) shows that \(B_1\le B_0\) at every belief, by Proposition~1(a). Apply it at the
belief \(M=M(\pi)\):
\begin{align*}
B_0(M)&\ge B_1(M) \why{Proposition~1(a) at the belief \(M\)}\\
&>c_H \why{right side of (H1)}.
\end{align*}
With the left side of (H1), \(B_0(\pi)<c_H<B_0(M)\). So \(B_0\) rises strictly between \(\pi\) and \(M\), and the
threshold \(\tau_0\) exists and lies between them. Against the weak incumbent, a price that lifted the belief to
\(M\) would recruit a high-cost challenger. No price does so, because the investor never trades (\(\Delta_0<k\)), and
every belief stays at \(\pi\). What the weak incumbent's economy lacks is information. Profit at the top posterior
would be enough there. The guide adds this observation; it follows from Proposition~1(a) and (H1).}

\answer{\ref{q:ccc-window}}{With \(\pi>1/2\), \(1-\pi<\pi\), so \(\min\{\pi,1-\pi\}=1-\pi\). By Step~1 of the
derivation in Section~\ref{sec:ccc-window},
\begin{align*}
\kappa(\pi)&=m(1-\pi) \why{\(\kappa(\pi)=m(\min\{\pi,1-\pi\})\)}\\
&=1-M(\pi) \why{\(1-M(\pi)=m(1-\pi)\)}.
\end{align*}
Also \(m(\pi)>m(1-\pi)\), because \(m\) rises strictly. The factor \(1-M(\pi)\) is the one in the lower bound
\(K_H\ge\rho[1-M(\pi)]\DT\), so class \(H\)'s purchase sets \(\kappa(\pi)\). In words, the market already leans toward
class \(H\). A purchase confirms the lean, so the doubt \(1-M(\pi)\) left after it is small. A sale goes against the
lean, and the doubt \(m(\pi)\) left after it is larger.}

\answer{\ref{q:ccc-silent}}{Take any candidate against \(F_0\), with any prices and any entry. The return of a purchase
is
\begin{align*}
K_H&=e(1)\,\Delta_0\,[1-\mu(1)] \why{trading return}\\
&\le\Delta_0 \why{\(e(1)\le1\), \(1-\mu(1)\le1\)}.
\end{align*}
So \(U_H(1)=K_H-k\le\Delta_0-k<0=U_H(0)\) by the left side of (H2). Class \(H\) prefers to abstain, and a candidate in
which it buys fails (E5). The expectations of the market makers act only through \(\mu(1)\), \(P(1)\) and \(e(1)\), and
the bound covers every value of these. The best case for the purchase gives class \(H\) the whole spread, and even
that is below \(k\). If market makers expect class \(H\) to buy, \(\mu(1)\) is high and \(1-\mu(1)\) is small. Even
with \(e(1)=1\), the return is at most \(\Delta_0[1-\mu(1)]\), which lies further below \(\Delta_0\).}

\answer{\ref{q:ccc-beta}}{The exclusion bound must hold under every investor strategy, and under some strategies the
posterior after the zero order exceeds \(M(\pi)\) unless \(b_0\ge b\), that is \(\beta\le1/3\) (Lemma~A.1(c)). Point 1
of Step~5 uses only the posteriors after nonzero orders (Lemma~A.1(b), any \(\beta<1/2\)), and it fixes full orders.
Then \(f_H(0)=f_L(0)=b_0\), so \(\mu(0)=\pi<M\) for every \(\beta\), all beliefs lie in \([m,M]\), and Step~0(f) gives
\(B_2<c_H\). Without point 1, point 3 would need a bound on \(\mu(0)\) under an arbitrary strategy. For
\(\beta>1/3\), the strategy in which class \(H\) abstains and class \(L\) never abstains gives \(\mu(0)>M(\pi)\) (the
proof of Lemma~A.1(c)).}

\answer{\ref{q:ccc-sufficient}}{Step~1 fails. With \(b+\lambda\le1/2\),
\[
e_H=\rho+(1-\rho)(b+\lambda)\le\rho+(1-\rho)\cdot\tfrac12=\frac{1+\rho}{2}.
\]
At the edge \(W_H(F_1)=2\rho W_H(F_0)/(1+\rho)\), with \(W_H(F_0)>0\),
\begin{align*}
e_H\,W_H(F_1)>\rho\,W_H(F_0)
&\iff e_H\cdot\frac{2\rho\,W_H(F_0)}{1+\rho}>\rho\,W_H(F_0) \why{substitute the edge}\\
&\iff e_H\cdot\frac{2}{1+\rho}>1 \why{divide by \(\rho W_H(F_0)>0\)}\\
&\iff e_H>\frac{1+\rho}{2} \why{multiply by \((1+\rho)/2>0\)}.
\end{align*}
Since \(e_H\le(1+\rho)/2\), the exact condition fails at the edge. So the first part of \eqref{eq:ccc-sufficient}
alone does not imply the exact condition, and \(b+\lambda>1/2\) does real work. The exact condition can still hold in
other cases. For instance, if \(W_H(F_1)=W_H(F_0)>0\), it reads \(e_H>\rho\), which always holds.}

\answer{\ref{q:ccc-decisive}}{Step~3 would break: with \(\bar\mu=1\) it reads \(c_H\le B(1)=g_H\), a weak
inequality. If \(g_H=c_H\), the high-cost challenger is indifferent, and under the default tie rule it prepares.
Decisive advantage speaks only when \(g_H>c_H\), so it would say nothing about \(W_H(F)\). In the model a liquidity
trader can place every order in both classes, so every belief is below one, and Step~3 is strict.}
